% concatenated from HWforCY.tex, notations.tex
\documentclass[11pt,leqno]{article}
\usepackage{graphicx, amsfonts, amsthm, amsxtra, amssymb, verbatim, makeidx}
\usepackage{subeqnarray, relsize}
\usepackage[mathscr]{euscript}
\usepackage{hyperref, tikz-cd}
\usepackage{aliascnt}
\usepackage{booktabs}
\usepackage{stmaryrd} %for the \llbracket and \rrbracket commands
\usepackage{soul} %for the \ul command. This gives you an underline, which runs over multiple lines
\usepackage[nameinlink,capitalize,noabbrev]{cleveref}
\hypersetup{
    colorlinks=true,       % false: boxed links; true: colored links
    linkcolor=blue,          % color of internal links
    citecolor=blue,        % color of links to bibliography
    filecolor=blue,      % color of file links
    urlcolor=blue           % color of external links
}


\textheight 24truecm
\textwidth 16truecm
\addtolength{\oddsidemargin}{-2truecm}
\addtolength{\evensidemargin}{-2truecm}
\addtolength{\topmargin}{-2truecm}
%----------------------------------------------------------------------
%-------------------------------------
\newtheorem{theo}{Theorem}
\newtheorem{coro}{Corollary}
\newtheorem{prob}{Problem}
\newtheorem{lemm}{Lemma}
\newtheorem{prop}{Proposition}
\newtheorem{conj}{Conjecture}
\newtheorem{exe}{Exercise}

\theoremstyle{remark}
\newtheorem{rem}{Remark}

\theoremstyle{definition}
\newtheorem{defi}{Definition}
\newtheorem{exam}{Example}

\makeindex
\makeglossary
\begin{document}




\def\U{{\cal U}}   %covering
\def\res{{\rm res}}
\def\cf{r}   %The coefficients of the linear cycles.


\def\Nn{{\sf n}}
\def\Am{{\sf A}}







%----------------General Math Notations-------------------------------------
\def\Z{\mathbb{Z}}                   %Integer  numbers
\def\Q{\mathbb{Q}}                   %Rational  numbers
\def\C{\mathbb{C}}                   %Complex numbers
\def\N{\mathbb{N}}                   %natural numbers
\def\Ff{\mathbb{F}}                  %finite field
\def\uhp{{\mathbb H}}                %upper half plane
\def\A{\mathbb{A}}                   %affine space C^n
\def\dR{{\rm dR}}                    %The subindex dR standing for de Rham cohomology.
\def\F{{\cal F}}                     %A foliation
\def\Sp{{\rm Sp}}                    %Symplectic group
\def\Gm{\mathbb{G}_m}                 %The multiplicative group
\def\Ga{\mathbb{G}_a}                 %The additive  group
\def\Tr{{\rm Tr}}                      %Trace map in Algebraic de Rham cohomology
\def\tr{{{\mathsf t}{\mathsf r}}}                 %Transposition of matrices
\def\spec{{\rm Spec}}            %The spectrume
\def\proj{{\rm Proj}}
\def\ker{{\rm ker}}              %kernel
\def\GL{{\rm GL}}                %The liner group


%----------------Gauss-Manin Connection in disguise---------------------
\def\k{{\sf k}}                     %Arbitrary field
\def\ring{{\sf R}}                   %A ring
\def\sk{{\mathfrak k }}             %Smaller field
\def\sring{{\mathfrak R }}          %Smaller ring


\def\X{{\sf X}}                      %families of varieties
\def\T{{\sf T}}                      %Moduli of enhanced  varieties
\def\V{{    V}}                   %Another parameter space

\def\Ts{{\sf S}}
\def\cmv{{\sf M}}                    %Classical moduli of varieties
\def\BG{{\sf G}}                       %Borel Algebraic Group
\def\podu{{\sf pd}}                   %Poincare dual
\def\ped{{\sf U}}                    %Period Domain
\def\per{{\sf  P}}                   %period matrix or fundamental system of GMC
\def\gm{{\sf  A}}                    %Gauss-Manin connection
\def\gma{{\sf  B}}                   %Gauss-Manin connection
\def\ben{{\sf b}}                    %Betti number

\def\Rav{{\mathfrak M }}                     % Space of modular vector fields
\def\Ram{{\cal C}}                         % Space of constant vector fields
\def\Rap{{i(\Lie(\BG))}}                    % Space of vector fields arising from group  action. 

\def\nov{{  n}}                    %Dimension of the variety
\def\mov{{  m}}                    %Fixed dimension of the cohomology
\def\Yuk{{\sf Y}}                     %Yukawa and generalizations 
\def\Ra{{\sf R}}                      %Ramanujan type vector field

\def\Da{{\sf D}}                      %Darboux vector field

\def\hn{{\sf h}}                      %Hodge numbers
\def\cpe{{\sf C}}                     %Constant periods
\def\g{{\sf g}}                       %An element of the Borel group
\def\t{{   t}}                       %An element \T
\def\v{{   v}}                       %An element \V


\def\pedo{{\sf  \Pi}}                  %Period domain before discrete group action


\def\Der{{\rm Der}}                   %Derivations
\def\MMF{{\sf MF}}                    %Moduli of modular foliations
\def\codim{{\rm codim}}                %codimension
\def\dim{{\rm    dim}}                %dimension
\def\Lie{{\rm Lie}}                   %Lie algebra of a group.
\def\gg{{\mathfrak  g}}                %elements of a Lie group


\def\u{{\sf u}}                       %An element of the period domain

\def\imh{{  \Psi}}                 %intersection matrix in homology
\def\imc{{  \Phi }}                  %intersection matrix in cohomology
\def\stab{{\rm Stab }}               %stablizer 
\def\Vec{{\Theta}}                 %space of vector fields
\def\prim{{0}}                  %primitive cohomology
\def\Zero{{\rm Zero}}                  %primitive cohomology


\def\Fg{{\sf F}}     %Genus g topological partition function
\def\hol{{\rm hol}}  %holomorphic
\def\non{{\rm non}}  %non-holomorphic
\def\alg{{\rm alg}}  %algebraic
\def\an{{\rm an}}   %analytic
\def\for{{\rm for}}  %holomorphic


\def\bcov{{\rm \O_\T}}       %The ring of modular-type functions

\def\leaves{{\cal L}}        %space of leaves  




\def\Hse{{\rm HS}}        %Hilbert series
\def\Hpo{{\rm HP}}        %Hilbert polynomial
\def\Hfu{{\rm HF}}        %Hilbert function
\def\Hsc{{\rm Hilb}}     %Hilbert scheme

\def\TS{\mathlarger{{\bf T}}}                %Tangent space
\def\IS{\mathlarger{{\cal I}}}                %Ideal sheaf 


\def\vf{{\sf v}}                      %A vector field
\def\wf{{\sf w}}                      %A vector field


\def\red{{\rm red}}                           %Gauss-Manin system


\def\Ua{{   L}}                      %An affine variety
\def\plc{{ Z_\infty}}    %a section of X by a projective space of dimension n/2+1








\def\gru{\mu} %the group of d-th roots of unity 
\def\pg{{ \sf S}}               %permutation group
\def\group{{ G}}            %Automorphisms of the Fermat



%---------------OLD Notation-------------------
\def\GM{{\rm GM}}


\def\perr{{\sf q}}        %period matrix.....
\def\perdo{{\cal K}}   %period domain
\def\sfl{{\mathrm F}} %Space of filtrations
\def\sp{{\mathbb S}}  %Sphere

\newcommand\diff[1]{\frac{d #1}{dz}} %Differential operator
\def\End{{\rm End}}              %Endomorphism group

\def\sing{{\rm Sing}}            %The set of singularities
\def\cha{{\rm char}}             %Charracteristic
\def\Gal{{\rm Gal}}              %The Galois group
\def\jacob{{\rm jacob}}          %the Jacobian ideal
\def\tjurina{{\rm tjurina}}      %the tjurina ideal
\newcommand\Pn[1]{\mathbb{P}^{#1}}   %Projective space of dimension #1
\def\P{\mathbb{P}}                   %Projective space
\def\Ff{\mathbb{F}}                  %Finite field



\def\O{{\cal O}}                     %ring of integers of a number field
\def\as{\mathbb{U}}                  %Some affine space
\def\ring{{\mathsf R}}                         %A ring
\def\R{\mathbb{R}}                   %real numbers



\newcommand\ep[1]{e^{\frac{2\pi i}{#1}}}% unipotent numbers
\newcommand\HH[2]{H^{#2}(#1)}        %Hodge structures
\def\Mat{{\rm Mat}}              %Matrices
\newcommand{\mat}[4]{
     \begin{bmatrix}
            #1 & #2 \\
            #3 & #4
       \end{bmatrix}
    }                                %two by two matrices
\newcommand{\matt}[2]{
     \begin{bmatrix}                 % one by two matrix
            #1   \\
            #2
       \end{bmatrix}
    }
\def\cl{{\rm cl}}                %Chern class

\def\hc{{\mathsf H}}                 %The set of Hodge cycles.
\def\Hb{{\cal H}}                    %Hodge bundle
\def\pese{{\sf P}}                  %Period set

\def\PP{\tilde{\cal P}}              %the period domain/ discrete group
\def\K{{\mathbb K}}                  %Field representing R or C


\def\M{{\cal M}}
\def\RR{{\cal R}}
\newcommand\Hi[1]{\mathbb{P}^{#1}_\infty}%the hyperplane at infinity
\def\pt{\mathbb{C}[t]}               %Polynomials in t
\def\W{{\cal W}}                     %weight filtration
%\def\Af{{\cal A}}                    %The space of tame polynomials.
\def\gr{{\rm Gr}}                %graded pieces
\def\Im{{\rm Im}}                %imaginary
\def\Re{{\rm Re}}                %Real
\def\depth{{\rm depth}}
\newcommand\SL[2]{{\rm SL}(#1, #2)}    %SL(2,Z)
\def\sl{{\rm SL}}                    %SL
\newcommand\PSL[2]{{\rm PSL}(#1, #2)}  %PSL(2,Z)
\def\Resi{{\rm Resi}}              %Residue

\def\L{{\cal L}}                     %The moduli of polarized lattices in a
                                     %fixed vector spaces.
\def\Aut{{\rm Aut}}              %Automorphism group of a vectorspace
\def\any{R}                          %Any subring of the field of complex
                                     %numbers.
\newcommand\ovl[1]{\overline{#1}}    %Conjugation of #1.



\newcommand\mf[2]{{M}^{#1}_{#2}}     %New modular functions
\newcommand\mfn[2]{{\tilde M}^{#1}_{#2}}     %New modular functions

\newcommand\bn[2]{\binom{#1}{#2}}    %Binomial
\def\ja{{\rm j}}                 %j of a two by two matrix
\def\Sc{\mathsf{S}}                  %Simple cycles
\newcommand\es[1]{g_{#1}}            %Eisenstein series
\newcommand\WW{{\mathsf W}}          %Similar to Milnor vector space
\newcommand\Ss{{\cal O}}             %Structural sheaf
\def\rank{{\rm rank}}                %rank of a module
\def\Dif{{\cal D}}                   %Differentials
\def\gcd{{\rm gcd}}                  %greatest common divisor
\def\zedi{{\rm ZD}}                  %zero divisors of a module
\def\BM{{\mathsf H}}                 %Brieskorn module
\def\plf{{\sf pl}}                             %Picard-Lefschetz formula
\def\sgn{{\rm sgn}}                      %sign
\def\diag{{\rm diag}}                   %diagonal matrix
\def\hodge{{\rm Hodge}}
\def\HF{{\sf F}}                                %The hodge filtration of the brieskon module
\def\WF{{\sf W}}                               %The weight filtration of the brieskon module
\def\HV{{\sf HV}}                                %humbert variety
\def\pol{{\rm pole}}                               %pole divisor
\def\bafi{{\sf r}}
\def\Id{{\rm Id}}                               %identity
\def\gms{{\sf M}}                           %Gauss-Manin system
\def\Iso{{\rm Iso}}                           %Gauss-Manin system






\def\hl{{\rm L}}    %holomorphic limit
\def\imF{{\rm F}}
\def\imG{{\rm G}}



\def\HL{{\rm Ho}}     %Hodge locus
\def\NLL{{\rm NL}}   %Noether-Lefschetz locus

\def\RG{{\bf G}}          %A Reductive group
\def\rg{{\bf g}}     %An element of a reductive group
\def\rbullet{{\cdot}}%Action of the reductive group.
\def\Ld{{\cal L}}      %Lie derivative 
%%%%%%%%%%%%%%%%%%%%%%%%%%%%%%%%My naming
\def\Ro{{\rm R}}     %Rogayeh
\def\ZS{{\rm ZeSc}}     %Zero scheme
\def\ZI{{\rm ZeId}}     %Zero ideal 
%\def\intint{{\mathlarger{\int}}}
 \def\integ{{\rm Int}}  %integral of an ideal

 
\def\tmap{{\sf t}}


\def\ivhs{{\rm IVHS}}    %IVHS ring
\def\ivhsmaps{{{\Delta}_{}}}   %The set of delta maps
\def\sch{{\rm Sch}}   %Scheme
\def\mk{{\mathfrak  m}}   %Maximal ideal
\def\pk{{\mathfrak  p}}   %prime ideal
\def\qk{{\mathfrak  q}}   %prime ideal




\newcommand\licy[1]{{\mathbb P}^{#1}} %linear projective cycle







\def\SS{\mathscr{S}}    %sheaf





\begin{center}
{\LARGE\bf Hasse-Witt invariants of Calabi-Yau varieties \footnote{\today}
}
\\
\vspace{.1in} {\large {\sc Jin Cao\footnote{School of Mathematics and Physics, University of Science and Technology Beijing, Beijing, China}, Mohamed Elmi\footnote{Yau Mathematical Sciences Center, Tsinghua University, Beijing, China} and Hossein Movasati}}
\footnote{
Instituto de Matem\'atica Pura e Aplicada, IMPA, Estrada Dona Castorina, 110, 22460-320, Rio de Janeiro, RJ, Brazil,
{\tt \href{http://w3.impa.br/~hossein/}{www.impa.br/$\sim$hossein}, hossein@impa.br.}}
\end{center}




%\tableofcontents



\begin{abstract}
We define the Hasse-Witt invariant of Calabi-Yau varieties in two different ways. The first method is through Cartier operator and the second method is through the theory of Calabi-Yau modular forms developed by the third author. We conjecture that these two definitions are equivalent and provide many examples of Calabi-Yau varieties in support of this conjecture. 
\end{abstract}

\section{Introduction}
Let $X$ be an elliptic curve over a perfect field $\sk$ of characteristic $p\not=2,3$ and let $\alpha$ be a regular differential $1$-form on $X$. The Cartier operator $C$ acts on the sheaf of differential $1$-forms on $X$ and the Hasse-Witt invariant is defined by the equality
$$
C(\alpha)={\rm HW}(X,\alpha)^{\frac{1}{p}}\alpha.
$$
The algebra of (elliptic) modular forms for $\SL 2\Z$ gives us another method for computing the Hasse-Witt invariant. It is well-known that this algebra is generated by the Eisenstein series $E_4$ and $E_6$. In particular, this implies that 
\begin{equation}
\label{19122025lastday}
E_{p-1}=A_p\left(\frac{1}{12}E_4,\frac{-1}{216}E_6\right)
\end{equation}
for some polynomial $A_p\in\Q[x,y]$ of weighted degree $p-1$ where $\deg(x)=4,\ \deg(y)=6$ . It turns out that $p$ does not appear in the denominators of the coefficients of $A_p$, we can reduce $A_p$ modulo $p$. If we write the pair $(X,\alpha)$ in the Weierstrass coordinates 
 \begin{equation}
 \label{25112025alirezaakbari}
 \X_{t_2,t_3}: y^2=4x^3-t_2x-t_3,\ \alpha=\frac{dx}{y}
 \end{equation} 
 then the following equality in characteristic $p$
 $$
{\rm HW(X,\alpha)}= A_p(t_2,t_3)
 $$
 is a classical result, mainly attributed to P. Deligne, see \cite[page 90]{ka73}. % and \cite[Theorem 3]{Serre1973}. 

The first part of the above story was developed by N. Katz for arbitrary varieties. In recent years, there has been some interest in similar notions in the case of Calabi-Yau varieties e.g. \cite{Ogus2001}. 

The third author has developed the theory of modular forms for Calabi-Yau varieties (CY modular forms) in many articles with the title "Gauss-Manin connection in disguise", see \cite{ho2020, GMCD-MQCY3} and the reference therein. Much of this work has been motivated by the study of generating series of Gromov-Witten invariants, even though there is no well-known enumerative algebraic geometry for the CY modular forms of the present article. 

The theory of CY modular forms generalizes the theory of quasi-modular forms $\Q[E_2,E_4,E_6]$. In this article, we deal only with part of the algebra of CY modular forms which generalizes $\Q[E_4,E_6]$.  It turns out that $\spec\left(\Z\left[\frac{1}{6(27t_3^2-t_2^3)}, t_2,t_3\right]\right)$ is the moduli scheme of pairs $(X,\alpha)$  and it gives us a geometric incarnation of the classical modular forms. 

We consider the moduli $\Ts$ of  pairs $(X,\alpha)$, where $X\subset\P^N_\C$ is a projective/polarized Calabi-Yau $n$-fold with a fixed Hilbert polynomial $P$, which varies in a component of the  Hilbert scheme ${\rm Hilb}_P(\P^N)$,  and $\alpha$ is a 
holomorphic differential $n$-form on $X$. 
The multiplicative group $\Gm:=(\C^*,\cdot)$ acts on $\Ts$ by:
 $$
(X,\alpha)\bullet a=(X,a^{-1}\alpha),\ a\in \Gm,\ (X,\alpha)\in \Ts.
$$
Let $L$ be the line bundle on the classical moduli space of Calabi-Yau $n$-folds given by holomorphic $n$-forms $\alpha$. The moduli space $\Ts$ is not the total bundle of $L$. It is such a total bundle quotiented by action of automorphisms $X$ on $\alpha$. We expect that 
\begin{conj}
\label{17102025bimsa-1}
The moduli space $\Ts$ has a natural affine scheme structure over $\Z[\frac{1}{N}]$ for some $N\in\N$. We have global regular functions $t_i$ on $\Ts$  and $k_i\in\Z$ such that
\begin{equation}
\label{02122025bri}
t_i(X,a^{-1}\alpha)=a^{k_i}t_i(X,\alpha), 
\end{equation}
the homogeneous polynomial $\Delta\in \Z[\frac{1}{N},t]:= \Z[\frac{1}{N},t_1,t_2,\cdots, t_s],\ \ \deg(t_i)=k_i$, called the discriminant,  such that 
$$
\Ts=\spec\left (\Z\left[ \frac{1}{N\Delta},t\right]/I\right ).
$$
Here, $I$ is a homogeneous ideal in $\Z[\frac{1}{N\Delta},t])$. Moreover, the moduli space is fine, that is, we have a universal family $\X\to\Ts$ which comes together with $\alpha\in H^0(\X,\Omega^n_{\X/\Ts})$. 
\end{conj}
We denote a fiber of $\X\to \Ts$ over $t\in\Ts$ by $\X_t$ which comes together with a unique differntial $n$-form $\alpha_t$. 
A global regular section of $\Ts$ is called a (meromorphic) CY modular form. 
E. Viehweg has constructed the coarse  moduli space of polarized smooth Calabi-Yau varieties defined over algebraically closed fields of characteristic zero as a quasi-projective variety, see \cite{viehweg95}. It does not seem to be difficult to manipulate his construction and prove \cref{17102025bimsa-1}. 
Note that once this conjecture is established, the classical moduli of Calabi-Yau varieties without the differential $n$-form $\alpha$ is the variety given by $I=0$ in the weighted projective space $\P^{k_1,k_2,\ldots, k_s}_\Q\backslash\{\Delta=0\}$. 

The above definition of CY modular forms does not indicate at all what $q$-expansion means. 
This is as follows. Let $X_z,\ z\in \C^m$ be a family of Calabi-Yau $n$-fold with $m:=h^{n-1,1}$ and assume that $z=0$ is a MUM point, see \cite{Morrison1993} for the definition.  Let $\alpha_z$ be a holomorphic $n$-form on $X_z$. We assume that all these are defined over $\Q$. It is well-known that we have cycles $\delta_{i,z}\in H_n(X_z,\Z),\ \ \ i=1,2,\ldots,m+1$ such that 
$$
\int_{\delta_{0,z}}\alpha_z,\ \ \int_{\delta_{i,z}}\alpha_z-\frac{\ln(z_i)}{2\pi\sqrt{-1}}\int_{\delta_{0,z}}\alpha_z,\ \ i=1,2,\ldots,m
$$ 
are holomorphic and the mirror map 
$$
(\C^m,0)\to(\C^m,0),\ \ \ (z_1,z_2,\ldots,z_m)\mapsto (q_1,q_2,\ldots, q_m),\ \ 
q_i:=\exp\left(\frac{\int_{\delta_{i,z}}\alpha_z}{\int_{\delta_{0,z}}\alpha_z}\right)
$$
is a biholomorphism, see \cite{Morrison1993}, and hence, we can invert it $z=z(q)$. We now take a smooth model of $X_z$ over $\Z[\frac{1}{N}]$ for some $N\in\N$, and  take $z$ from a perfect field $\sk$ of characteristic $p\nmid N$ and take $(X_z, \alpha_z)$ defined over $\sk$. We have the Cartier operator $C$ acting on closed differential forms, and $C(\alpha_z)$ is also a global holomorphic $n$-form on $X_z$. Therefore, 
$$
C(X_z, \alpha_z)={\rm HW}(X_z,\alpha_z)^\frac{1}{p}\alpha_z,\ \ \ a(z):={\rm HW}(X_z,\alpha_z)
$$
where ${\rm HW}(X_z,\alpha_z)\in \sk$ is called the Hasse-Witt invariant. It turns out that in all our examples  the Hasse-Witt invariant is a polynomial  in $z$ with coefficients in $\Ff_p$. From another side we have the period integral 
$$
{\rm HP}(X_z,\alpha_z):=
\int_{\delta_{0,z}}\alpha_z\Big / \int_{\delta_{0,0}}\alpha_0
$$
which is the unique holomorphic solution to the Picard-Fuchs system of $(X_z,\alpha_z)$ with the constant term $1$, and moreover, it is $p$-integral, and hence we can take its reduction modulo $p$. It turns out that both quantities ${\rm HP}(X_z,\alpha_z)$ and ${\rm HW}(X_z,\alpha_z)$ can be evaluated at $z=0$.


% and in the following conjecture we need their normalized versions with valuation at $z=0$ equal to $1$.   
%\marginpar{\tiny In order to make the following conjecture true, we need to make a modification of the above ${\rm HP}(X_z,\alpha_z)$. The reason is as follows. Apply the above construction to Weierstrass family $y^2=4x^3-t_2x-t_3$ and view the Legrend family as one parameter family. Then $t_3 = \frac{4}{27}(2z^3-3z^2-3z+2)$. We already know that $t_3 = (-\frac{1}{216})E_6$. On the other hand, we should also have an equality $t_3 = {\rm HP}^6[\frac{4}{27}(2z^3-3z^2-3z+2)]$. The latter formula become true only if $HP = \frac{i}{2} F(\frac{1}{2},\frac{1}{2}, 1;z)$. In other words, if $HP$ is with the constant term $1$, we could never get the $-1$ in the expression $t_3 = (-\frac{1}{216})E_6$.}
\begin{conj}\rm
\label{17112025ymsc}
For all primes $p\not| N$,
the Hasse-Witt invariant ${\rm HW}(X_z,\alpha_z)$ up to sign  is the truncation of the holomorphic period at degree $p-1$  and  the quantity
$$
A_p(z):={\rm HP}(X_z,\alpha_z)^{p-1}
{\rm HW}(X_z,\alpha_z)
$$
after inserting the mirror map $A_p(z(q))$ has $p$-integral coefficients and modulo $p$ it is a  constant independent of $z$, more precisely, $A_p(z(q))\equiv_p 1$ (up to sign).
%\marginpar{\tiny Still not happy, it might be  $-1$} 
\end{conj}
For the first part of Conjecture \ref{17112025ymsc}, it has been shown in the case of Calabi-Yau varieties realized as some hypersurfaces in a smooth toric variety in \cite{HLYY23} and some three-dimensional reflexive polytopes (see \cite[Main Theorem]{SalernoWhitcher22}). However the $p$-integrality of the mirror map itself, which implies the $p$-integrality of $A_p(z(q))$, is still an unsolved conjecture which has been settled mainly in the case of hypergeometric Calabi-Yau varieties using results of B. Dwork, see \cite[Appendix C]{GMCD-MQCY3} and references therein. This is going to be discussed in the third author's paper \cite{Ibiporanga} in which a larger moduli space (ibiporanga) as in \cite{ho2020, GMCD-MQCY3} is used.    

The first part of the above conjecture might not be so difficult as it only requires the study of the degeneration of Hasse-Witt invariant when $X_z$ becomes singular, see \cref{27112025bimsa}. However, note that in \cref{17112025ymsc} we cannot replace the holomorphic period with its truncated one, and hence consider $A_p(z)=a(z)^{p}$,  since the $q$-expansion has to do with the full period.


Let us now rewrite \cref{17112025ymsc} in the context of \cref{17102025bimsa-1}. Since $\Ts$ is the moduli space of the pairs $(X,\alpha)$, for each $z$ in a small neighborhood of $0\in\C^n$ with $z_i\not=0$, we have a unique $t(z):=(t_1(z),t_2(z),\ldots,t_s(z))\in\Ts$, where $t_i(z)$'s are rational functions in $z$,  such that 
$$
(X_z,\alpha_z)\cong (\X_{t(z)},\alpha_{t(z)}).
$$
We define a special locus
$$
\tilde\Ts:=\left\{ t\in \Ts \ \ \ \ \Big| \ \ \ \ 
\int_{\delta}\alpha_t
=c,\ \ \hbox{for some cycle} \ \  \delta\in H_n(\X_t,\Z) \right \}.
$$
where $c:=\int_{\delta_{0,0}}\alpha_0$ as above. It plays the similar role as the $\tau$-locus defined in \cite[Section 4.2]{GMCD-MQCY3}.
Note that $ \int_{\delta}\alpha_t$ is a holomorphic multi-valued function in $\Ts$. 
By the functional equation of $t_i$'s in \eqref{02122025bri} it follows that such a loci in $t_i$-coordinates is given by %\marginpar{\tinyShall we explain why introduce $\tilde\Ts$?}
$$
\left ({\rm HP}(X_z,\alpha_z)^{k_i}t_i(z),\ \ i=1,2,\ldots,s\right )\in\Ts.
$$
After inserting the mirror map we get $q$-expansions of $t_i$'s that we denote by $t_i(q)$. This is the holomorphic incarnation of the algebraic CY modular forms $t_i$ as a regular function on $\Ts$. \cref{17112025ymsc} turns into the following: 
\begin{conj} \label{02112025huairou}
For any prime $p\nmid N$ we have a homogeneous polynomial $A\in \Ff_p[t]$ of degree $p-1$ with $\deg(t_i)=k_i$ such that $C(\alpha)=A_p(t)^{\frac{1}{p}}\alpha$, where $C$ is the Cartier operator. Moreover, if we write the $q$-expansion of $t_i$'s then 
$$
A_p(t_1(q),t_2(q),\cdots, t_s(q))\equiv_p1. 
$$
\end{conj}
We recall that one way to prove \cref{02112025huairou} for elliptic curves is to observe that $A_p(t_2,t_3)$ is a Hecke eigenform of weight $p-1$ with the constant term $1$. This determines $A_p$ uniquely. Unfortunately, all the efforts of the third author to generalize Hecke operators for Calabi-Yau modular forms has failed, see \cite[Section 11.2]{GMCD-MQCY3}. %\marginpar{\tiny \color{red} Jin. This is an easier proof.}
We are able to automate a verification of our conjetures and find that
% \begin{theo}
% \label{28112025bri}
% %Let $M=200$ and $N=200$. 
% We have the following statements valid for the first $200$ primes and truncation of $q$-series at $q^{200}$. %\marginpar{\tiny  {\color{red} Mohamed} We have to put explicit values of $N$ and $M$ and run our code in a better computer!!}
% \begin{enumerate}
% \item 
% \cref{17102025bimsa-1}, \cref{17112025ymsc} and \cref{02112025huairou} are true for four families of hypergeometric Calabi-Yau threefolds in  
% \cite{Morrison:1991cd}. We have $\Ts_\Q=\spec(\Q[t_1, t_k, \frac{1}{k(t_1^k-t_k)}])$. The universal family of hypersurfaces in the weighted projective space $\mathbb{P}^{k_0,k_1,k_2,k_3,k_4}$ is given in table below:

% \begin{equation*}
% \begin{array}{|c|c|c|}
% \hline
% \cite{alenstzu} \# & (k_0,k_1,k_2,k_3,k_4) &  f\\[2pt]
%    \hline
% 1 & (1,1,1,1,1) & t_5x_0^5 + x_1^5 + x_2^5+x_3^5 + x_4^5 -t_1x_0x_1x_2x_3x_4\\[2pt]
% 2 & (5,2,1,1,1) & x_0^2 + x_1^5 + t_{10}x_2^{10} + x_3^{10} + x_4^{10} - t_1 x_0x_1x_2x_3x_4\\[2pt]
% 7 & (4,1,1,1,1) & x_0^2 + t_8x_1^8 + x_2^8 + x_3^8 + x_4^8 - t_1x_0x_1x_2x_3x_4\\[2pt]
% 8 &   (1,1,1,2,1) & t_6x_0^6 + x_1^6 + x_2^6 + x_3^3 + x_4^6 - t_1x_0x_1x_2x_3x_4\\[2pt]
% \hline
% \end{array}
% \end{equation*}
% The differential $3$-form is given by 
% $\alpha=\frac{dx_1\wedge dx_2\wedge dx_3\wedge dx_4}{df}$, 
% where $f$ is the equation of the Calabi-Yau threefold  with $x_0=1$. 
% \item 
% \cref{17112025ymsc} is true for all the Calabi-Yau equations by AESZ in \cite{alenstzu}. 
% \end{enumerate}
% \end{theo}
\begin{theo}
\label{28112025bri}
%Let $M=200$ and $N=200$. 

\cref{17102025bimsa-1}, \cref{17112025ymsc} and \cref{02112025huairou} are true for the first $200$ primes and to $O(q^{200})$ for four families of hypergeometric Calabi-Yau threefolds in  
\cite{Morrison:1991cd}. We have $\Ts_\Q=\spec(\Q[t_1, t_k, \frac{1}{k(t_1^k-t_k)}])$. The universal family of hypersurfaces in the weighted projective space $\mathbb{P}^{k_0,k_1,k_2,k_3,k_4}$ is given in table below:

\begin{equation*}
\begin{array}{|c|c|c|}
\hline
\cite{alenstzu} \# & (k_0,k_1,k_2,k_3,k_4) &  f\\[2pt]
   \hline
1 & (1,1,1,1,1) & t_5x_0^5 + x_1^5 + x_2^5+x_3^5 + x_4^5 -t_1x_0x_1x_2x_3x_4\\[2pt]
2 & (5,2,1,1,1) & x_0^2 + x_1^5 + t_{10}x_2^{10} + x_3^{10} + x_4^{10} - t_1 x_0x_1x_2x_3x_4\\[2pt]
7 & (4,1,1,1,1) & x_0^2 + t_8x_1^8 + x_2^8 + x_3^8 + x_4^8 - t_1x_0x_1x_2x_3x_4\\[2pt]
8 &   (1,1,1,2,1) & t_6x_0^6 + x_1^6 + x_2^6 + x_3^3 + x_4^6 - t_1x_0x_1x_2x_3x_4\\[2pt]
\hline
\end{array}
\end{equation*}
The differential $3$-form is given by 
$\alpha=\frac{dx_1\wedge dx_2\wedge dx_3\wedge dx_4}{df}$, 
where $f$ is the equation of the Calabi-Yau threefold  with $x_0=1$. 

\end{theo}



% $$
% \begin{array}{c|c|c}
%    AESZ \# & \mathbb{P}[k_i] &  P(x,\psi)\\
%    \hline
% 1 & (1,1,1,1,1) & M_5x_0^5 + x_1^5 + x_2^5+x_3^5 + x_4^5 -M_1x_0x_1x_2x_3x_4\\
% 2 & (1,1,1,2,5) & M_{10}x_0^2 + x_1^5 + x_2^{10} + x_3^{10} + x_4^{10} - M_1 x_0x_1x_2x_3x_4\\
% 7 & (1,1,1,1,4) & M_8x_0^2 + x_1^8 + x_2^8 + x_3^8 + x_4^8 - M_1x_0x_1x_2x_3x_4\\
% 8 &   (1,1,1,1,2) & M_6x_0^6 + x_1^6 + x_2^6 + x_3^3 + x_4^6 - M_1x_0x_1x_2x_3x_4\\
% \end{array}
% $$

%As in this article we are interested in explicit construction of $\Ts$,  we will verify \cref{17102025bimsa-1} and the following conjecture for four examples of hypergeometric Calabi-Yau threefold. 
%We can formulate \cref{02112025huairou} without the construction of  the moduli space $\Ts$. 




We return to the case of elliptic curves and explain \cref{17112025ymsc} for the Legendre family $E_z:y^2=x(x-1)(x-z)$ together with the holomorphic 1-form $\alpha_z = \frac{dx}{y}$. The $1$-form $\alpha_z$ can be integrated over generators $\delta_{0,z}$ and $\delta_{1,z}$ of $H_1(E_z,\mathbb{Z})$. This leads to the  $A$ and $B$ periods
%\marginpar{\tiny back to $\delta$ notation and not A,B}
% \begin{eqnarray*}
% \int_A\alpha_z&=2\int^{\infty}_1 \alpha_z 
% &= 2\pi F(
% \frac{1}{2} \frac{1}{2},1; z) =
% 2\pi\sum_{k=0}^\infty \frac{(\frac{1}{2})_k (\frac{1}{2})_k}{(k!)^2} z^k\\
% %\pi(\sum_{i=0}^{\infty} \frac{1}{2^{4i}}\binom{2i}{i}^2 z^i)
% \\
% \int_B\alpha_z &=\sqrt{-1}\int^{0}_{-\infty} \alpha_z &= F(\frac{1}{2} \frac{1}{2},1; z)\ln(z)+G(z),\ 
%  \ G(z)=\sum_{k=1}^\infty \frac{(\frac{1}{2})_k (\frac{1}{2})_k}{(k!)^2}\big{[}\sum_{j=1}^ 2\sum_{i=0}^{k-1}(\frac{1}{\frac{1}{2}+i}-\frac{1}{1+i})\big{]}z^k.
% \end{eqnarray*}
\begin{align*}
\int_{\delta_{0,z}}\alpha_z &= 2\pi F\left(\frac{1}{2} \frac{1}{2},1; z\right)=
2\pi\sum_{k=0}^\infty \frac{(\frac{1}{2})_k (\frac{1}{2})_k}{(k!)^2} z^k
%\pi(\sum_{i=0}^{\infty} \frac{1}{2^{4i}}\binom{2i}{i}^2 z^i)
\\
\int_{\delta_{1,z}} \alpha_z &= \frac{2}{\sqrt{-1}}\left(F\left(\frac{1}{2} \frac{1}{2},1; z\right)\ln\left(\frac{z}{16}\right)+G(z)\right),\ \ \mathbb{ }G(z)=4\sum_{k=1}^\infty 
 \frac{(\frac{1}{2})_k (\frac{1}{2})_k}{(k!)^2}
(H_{2k}-H_k)z^k.
\end{align*}
where $H_k$ is the $k^\text{th}$ Harmonic number. See, for example, \cite{Zagier2018ArithmeticTopologyDE}. The periods are annihilated by the Picard-Fuchs operator
\[
%\mathcal{D} =
%z(1-z) \frac{d^2}{d z^2} + (1-2z) \frac{d}{dz}-\frac{1}{4}=
\theta^2-z\left(\theta+\frac{1}{2}\right)^2=0,\ \ \ \ \theta:=z\partial_z
~\]
and, by construction, they have monodromy in $SL(2,\mathbb{Z})$. 

Let $p$ be a prime bigger than $3$. The Hasse-Witt invariant of the Legendre family is simply the truncation of the holomorphic period: 
\[
a(z) = (-1)^{\frac{p-1}{2}}
\sum_{k=0}^\frac{p-1}{2} \frac{(\frac{1}{2})_k (\frac{1}{2})_k}{(k!)^2} z^k,
%(-1)^{\frac{p-1}{2}}\sum^{\frac{p-1}{2}}_{i=0} \binom{\frac{p-1}{2}}{i}^2 z^i.
\]
which, modulo $p$, is the coefficient of $x^{p-1}$ in the expression $(x(x-1)(x-z))^{\frac{p-1}{2}}$. 

We define $\tau(z)$ as the ratio of $B$ and $A$ periods and compute the mirror map $q(z)=\frac{z}{16}\exp\left(\dfrac{G(z)}{F(z)} \right)
$.\footnote{Note that $q=e^{\pi i \tau}$.} As expected, we recover the $q$-expansion of the modular lambda function when we invert this power series i.e.
\begin{equation*}
    z(q)=\frac{\theta^4_2(q)}{\theta^4_3(q)}=16 q-128 q^2+704 q^3-3072 q^4+11488 q^5-38400
   q^6+O\left(q^7\right),
\end{equation*}
We verify that, for the first $200$ primes,
$$
A_p(z(q)) = F\left(\frac{1}{2} \frac{1}{2},1; z(q)\right)^{p-1}a(z(q))=(-1)^{\frac{p-1}{2}} + \mathcal{O}(q^{200})  \mod p.
$$
If we consider the moduli space $\Ts$ of $(E,\alpha)$ as before, the previous procedure gives us 
$$
%\frac{4}{3}
\left(z(q)^2-z(q)+1\right) F\left(\frac{1}{2} \frac{1}{2},1; z(q)\right)^4=E_4(q^2),\ 
$$
$$
%\frac{8}{27}
\left(z(q)^3-\frac{3}{2}z(q)^2-\frac{3}{2}z(q)+1\right)F\left(\frac{1}{2} \frac{1}{2},1; z(q)\right)^6=E_6(q^2).
$$
%From these equalities and $A(z(q))\equiv_p 1$ we get  \cref{19122025lastday}. Note that $(-1)^{\frac{p-1}{2}}$ is absolved in the definition of $A$ is 
\begin{rem}
It is known that:
$$
F\left(\frac{1}{2} \frac{1}{2},1; z(q)\right) = \theta_3^2(q), \ \ \ \theta_3(q) = \sum^{\infty}_{k=-\infty} q^{k^2}.
$$
For example, see \cite[(4.32)]{Wenzhe2021} or \cite[Theorem 8.3]{Chan2020}.
%Combining with (7.28) and (7.29) in \cite{Chan2020},
%\[
%E_4(q^2) = (z(q)^2-z(q)+1) \theta_3^8(0|q)
%\]
%and
%\[
%E_6(q^2) = (z(q)^3-\frac{3}{2}z(q)^2-\frac{3}{2}z(q)+1) \theta_3^{12}(0|q)
%\]
%we could also prove the above formulas.
\end{rem}
%(z(q)^2-z(q)+1) F\left(\frac{1}{2} \frac{1}{2},1; z(q)\right)^4 = (z(q)^2-z(q)+1) \theta_3^8(0|q) = E_4(q^2)
%$$
%and 
%$$
%(z(q)^3-\frac{3}{2}z(q)^2-\frac{3}{2}z(q)+1)F\left(\frac{1}{2} \frac{1}{2},1; z(q)\right)^6 = (z(q)^3-\frac{3}{2}z(q)^2-\frac{3}{2}z(q)+1) \theta_3^{12}(0|q) = E_6(q^2).
%$$
In a similar way, if we consider the moduli of $(E,\alpha,P,Q)$, where $(E,\alpha)$ as before, and $P$ and $Q$ are 2-torsions with Weil pairing $+1$ then the moduli space is $\spec(\Z[\frac{1}{6t_2s_2(t_2-s_2)},t_2,s_2])$ with the universal family $y^2=x(x-t_2)(x-s_2)$ and we get 
$$
\frac{1}{4}F\left(\frac{1}{2} \frac{1}{2},1; z(q)\right)^2:=q\frac{\partial}{\partial q}\ln(\frac{\theta_2(0|q)}{\theta_4(0|q)})
$$
and
$$
\frac{z(q)}{4}F\left(\frac{1}{2} \frac{1}{2},1; z(q)\right)^2=q\frac{\partial}{\partial q}\ln(\frac{\theta_3(0|q)}{\theta_4(0|q)}),
$$
which generate the algebra of  modular forms  for $\Gamma(2)$, see \cite[Page 334]{ho14}. 

\begin{comment}
ring r=(0,z),x,dp;
4*(x+(z+1)/3)*(x+(z+1)/3-1)*(x+(z+1)/3-z);
\end{comment}
%\begin{rem}
%We also have:
%$$
%\frac{1-z(q)}{6}F\left(\frac{1}{2} \frac{1}{2},1; z(q)\right)^2:=q\frac{\partial}{\partial q}\ln(\frac{\theta_3(0|q)}{\theta_2(0|q)}).
%$$
%\end{rem}


%We have 
%\[
%(-1)^{\frac{p-1}{2}} a(z)=\sum_{i=0}^{\frac{p-1}{2}} \frac{1}{2^{4i}}\binom{2i}{i}^2 z^i, \mod p. \]
%the latter operator also annihilates $a(z)$ modulo $p$, see the proof of \cite[Theorem 4.1]{si09}, see also \cite[Proposition 2.3.6.3]{Katz1972}. 




{\bf Acknowledgment:} This article was written during the third author's visit to BIMSA and YMSC at Tsinghua University. We thank S.T.Yau for his efforts in creating these insitutions, where we could meet and discuss mathematics. The first author was supported by the University of Science and Technology Beijing Foundation, China (Grant No. 00007886) and the Fundamental Research Funds for the Central Universities (Grant No. 06320202). The second author would like to acknowledge support from the Shuimu Tsinghua Scholar Program and the China Postdoctoral Science Foundation.

\section{Cartier operator and Hasse-Witt invariant}
\label{27112025bimsa}
In this section we review Cartier operator and Hasse-Witt invariant from \cite{BK05}.
Let $X = \mathrm{Spec}(A)$ be an affine scheme over a perfect field $\sk$ of characteristic $p >0$ and $\Omega_A^1$ be the module of K\"{a}hler differentials of $A$ over $k$, which is equipped with the $\sk$-derivation
\[
d: A \to \Omega^1_A, a \to da.
\]
Then the de Rham complex $(\Omega_A^{\bullet}, d)$ of $A$ consists of the exterior algebra $\Omega_A^{\bullet}$ of $\Omega_A^1$ over $A$ together with the extended maps $d: \Omega^{\bullet}_A \to \Omega^{\bullet}_A$. 
We let
\[
Z^i_A = \{\alpha \in \Omega_A^i \mid d\alpha = 0\},  \ \ \ B^i_A = d \Omega_A^{i-1}, H_A^i = Z^i_A/B^i_A.
\]
Then one may define an $A$-algebra homomorphism:
\begin{eqnarray*} \nonumber
    \gamma: \Omega^{\bullet}_A &\to& H^{\bullet}_A \\
    a_1da_2 \wedge \cdots\wedge da_n &\to& a_1^p a_2^{p-1}da_2 \wedge \cdots \wedge a_n^{p-1} da_n (\mathrm{mod}\ B_A^{\bullet}), 
\end{eqnarray*}
where the $A$-module structure on $H^{\bullet}_A$ comes from the Frobenius action on $A$, see \cite[Lemma 1.3.3]{BK05}. Now we consider $X$ is a scheme over $\sk$ and let $F$ be the relative Frobenius $X \to X^{(p)}$. Then the sheaf-theoretic version of the above construction provide us a unique homomorphism of sheaves of graded-commutative $\mathcal{O}_{X^{(p)}}$-algebras:
\[
\gamma: \Omega^{\bullet}_{X^{(p)}} \to \bigoplus^{\infty}_{i=0} \mathcal{H}^i F_{*} \Omega_X^{\bullet}
\]
which satisfies:
\begin{itemize}
    \item 
    $\gamma(f) = f^p$ for $f \in \mathcal{O}_{X^{(p)}}$;
    \item 
    $\gamma(df) = f^{p-1}df(\mathrm{mod} \ d\mathcal{O}_{X^{(p)}})$ for $f \in \mathcal{O}_{X^{(p)}}$;
    \item 
    For $\alpha, \beta \in \Omega^{\bullet}_{X^{(p)}}$, we have: $\gamma(\alpha + \beta) = \gamma(\alpha) + \gamma(\beta)$ and
    $\gamma(\alpha \wedge \beta) = \gamma(\alpha) \wedge \gamma(\beta)$.
\end{itemize}
We remark that $\gamma$ is an isomorphism when $X$ is a nonsingular variety. The proof is contained in \cite[Theorem 2.1.1]{Katz1972} or \cite[Theorem 1.3.4]{BK05}.
\begin{defi}
Let $X$ is a nonsingular variety of dimension $n$. The inverse of the isomorphism $\gamma$ is called the Cartier operator:
\[
C = \sum^n_{i=0} C_i: \bigoplus^{n}_{i=0} \mathcal{H}^i F_{*} \Omega_X^{\bullet} \to \Omega^{\bullet}_{X^{(p)}}.
\]
\end{defi}
\begin{rem}
For $X$ a nonsingular variety of dimension $n$, we define the Cartier operator for $n$ forms as the composition of the maps $C: F_* \Omega_X^n \to \mathcal{H}^n F_{*} \Omega_X^{\bullet} \to \Omega^{n}_{X^{(p)}}$ by abuse of notations.
\end{rem}

We let $X$ be a smooth family of nonsigular projective Calabi-Yau varieties of dimension $n$ over an integral affine variety $S = \mathrm{Spec}(R)$, where $R$ is a finite generated algebra over a perfect field $\sk$ of characteristic $p > 0$. Then $f_* \Omega_{X/S}$ is a locally free sheaf of rank $1$ over $S$, where $f: X \to S$. Now we fix a basis $\alpha$ of $f_* \Omega_{X/S}$, whose dual basis under the Serre duality is denoted by $\eta \in R^nf_* \mathcal{O}_X$. Note that the absolute Frobenius action $F_{\mathrm{abs}}: \mathcal{O}_X \to \mathcal{O}_X$ induces an endomorphism $F_{\mathrm{abs}}^*$ of $R^nf_* \mathcal{O}_X$. Then:
\begin{defi}
We define the Hasse-Witt invariant of $X$ with respect to $\alpha$ as the element $A_p(X, \alpha) \in R$ such that:
\[
F_{\mathrm{abs}}^*(\eta) = A_p(X, \alpha) \eta.
\]
\end{defi}
\begin{rem}
If we assume further that each fiber of $X$ is a hypersurface, then the Hasse-Witt invariant can be also defined as 
\[
C(F_* \omega) =  A_p(X, \alpha)^{\frac{1}{p}} \omega^{(p)}.
\]
The proof of this equality is based on the algorithm of the computation of Hasse-Witt invariants for hypersurfaces of Katz \cite[Algorithm 2.3.7.14]{Katz1972} and \cite[Corollary 1]{Miller72}, since both of these computational results lead to the same invariant.
\end{rem}

One remarkable property is that the Hasse-Witt invariant satisfies the Picard-Fuchs equation due to the Katz-Igusa-Manin theorem \cite[Proposition 2.3.6.3]{Katz1972}. Applying this result to the family of Calabi-Yau $n$-folds, we have:
\begin{theo}
\label{01122025zutopia}
Let $\mathcal{D}$ be $\sk$-linear differential operator on $S$, which is contained in the algebra generated by $\mathrm{Der}(S/\sk)$ and acts on the de Rham cohomology sheaves $f_* \Omega_{X/S}$ via the Gauss-Manin connection $\nabla$. Suppose that
\[
\nabla_{\mathcal{D}}(\alpha) = 0,
\]
where we view $\alpha$ as a section of $f_* \Omega_{X/S}$ as above.
Then $\mathcal{D}(A_p(X, \alpha)) = 0$.
\end{theo}
Note if ${\cal D}:= \sum p_i\frac{\partial^i}{\partial t^i}$ then by definition we have $\nabla_{\cal  D}=\sum p_i\nabla_{\frac{\partial}{\partial t}}^i$. From the above theorem it might be possible to prove the following: 
 The Hasse-Witt invariant of $X_z$ is up to multiplication by a constant  $c_2\in\Ff_p$ 
 the truncation at order $p-1$ of the Taylor sereis of the holomorphic period $\psi_0(z):=(2\pi i)^{-n}\int_{\delta_{0,z}}\alpha_{z}$ at the MUM point. The idea of the proof is as follows: 
 Let us consider the differential sub module $L$ of $\Z[\frac{1}{N}][z_1,z_2,\ldots,\theta_1,\theta_2,\ldots]$ (or Picard-Fuchs system) which annihilates the period $\psi_0(z)$. 
By \cref{01122025zutopia} we know that the Hasse-Witt invariant $a(z)$ satisfies also $L(a(z))=0$. A priori, $a(z)$ might have poles along $z_i=0$, and other degeneracy loci of $X_z$. We must prove that this is not the case, and hence $a(z)$ is a polynomial in $z$. %\marginpar{\tiny Comment: If we consider $A$ with the coordinates in the $\sf T$-space, the conjecture 1 together with the definition of the Hasse-Witt invariant would imply $A(t_i) \in \mathcal{O}_S$, which is a polynomial.}  
Since $0$ is a MUM point, we know that there is a unique meromorphic solution $\psi_0(z)$ to $L=0$.
%This is also true considering modulo $p$-reduction of $L$. This together with the fact that $a(z)$ is a polynomial in $z$  implies that $a(z)\equiv_p \psi_0(z)$, 
From all these we might try to get the desired statement. We  must further argue that this polynomial has degree $\leq p-1$.   
%Two MUM points
%Supersingular CY3 and Ogus article! Definition of Cartier operator for $n$ forms on $n$-dimensional varieties. 
 %Definition of Hasse-Witt invariant for Calabi-Yau $n$-folds. 
 %Katz-Igusa-Manin theorem on PF and HW. \marginpar{\tiny{\color{red} Jin} write about these topics as short and readable as possible}












\section{Mirror quintic}
Recall that in \cite[Section 3]{GMCD-MQCY3}, the third author constructed the moduli $\sf S$ of pairs $(X, \alpha)$, where $X$ is a mirror quintic and $\alpha$ is a holomorphic differential 3-form on $X$. He further showed that
\[
\Ts_\Q \cong \mathrm{Spec}(\mathbb{Q}[t_1, t_5, \frac{1}{(t_1^5-t_5)t_5}])
\]
and the universal family over $\Ts_\Q$ is given by a desingularization of $X_t: \P\{f=0\}/G$, where
\[
f(x) = -t_5x_0^5 - x_1^5 - x_2^5 - x_3^5 - x_4^5+5t_1x_0x_1x_2x_3x_4
\]
and $\alpha$ in the affine coordinate $x_4=1$ is given by
\[
\alpha= \frac{dx_0 \wedge dx_1 \wedge dx_2} 
{\partial f/\partial x_3}.
\]
The holomorphic period of $\alpha$ along the 3-torus \[\delta_{0,z} = \{(x_0, x_1, x_2, x_3, x_4) \mid |x_0| = |x_1| = |x_2| = \delta, |x_3| \ll 1, x_4=1\}\] for $t_1=1, t_5=z$ 
is given by
$$
5 (2\pi i)^{-3} \int_{\delta_{0,z}} \alpha_z = \sum^{\infty}_{k=0} \frac{(5k)!}{(k!)^5} (\frac{z}{5^5})^k, \quad |z|<1.
$$

More generally, we may consider an arbitrary Dwork family of n-folds in $\mathbb{P}^{n+1}$, together with a choice of holomorphic $n$-form $\alpha$.\cite{GMCD-DF} This leads us to the family 
\[
X_{t_0,t_{n+2}}:f=0,\ \  
f=-t_{n+2}x_0^{n+2} - x_{1}^{n+2}  - \cdots - x_{n+1}^{n+2} + (n+2)t_1x_0x_1\cdots x_{n+1}.
\] 
\[ 
\alpha= \frac{dx_0 \wedge \cdots \wedge dx_{n}}{f_{n}}, f_{n} = \frac{\partial f}{\partial x_{n}},\  \ \hbox{ in the affine coordinate } x_{n+1}=1
\]
The holomorphic period of $\alpha$ for $t_1=1, t_{n+2}=z$  is given by
%$$
%F(z):=(\frac{2\pi i}{n+2})^{-n-1}\int_{\delta_{0,z}}\alpha_z=\sum_{k=0}^\infty \frac{((n+2)k)!}{k!^{n+2}} \, z^k, \quad |z|<1
%$$
%The formula is 
$$
F(z) = (n+2) (2\pi i)^{-n} \int_{\delta_{0,z}} \alpha_z = \sum^{\infty}_{k=0} \frac{((n+2)k)!}{(k!)^{n+2}} (\frac{z}{(n+2)^{n+2}})^k, \quad |z|<1.
$$




%\marginpar{\tiny {\color{red} Jin} We have to check $2\pi i$ factor, see the article \cite{GMPR}.  }

\begin{prop} We let $p$ prime  $p\not|n+2$. Then:
\begin{equation}
C(\alpha) = \left( \sum_{k=0}^{\lfloor \frac{p-1}{n+2} \rfloor} \frac{((n+2)k)!}{(k!)^{n+2}}
%\binom{p-1}{(n+2)k} 
t_{n+2}^k ((n+2)t_1)^{p-1-(n+2)k} \right)^{\frac{1}{p}} \alpha
\end{equation}
\end{prop}

%\begin{prop}
%We have:
%\begin{equation}
%C(\omega) = \left( \sum_{i=0}^{\lfloor \frac{p-1}{5} \rfloor} \frac{(5i)!}{(i!)^5}\binom{p-1}{5i} (-t_4)^i(5t_0)^{p-1-5i} \right)^{\frac{1}{p}} \omega,
%\end{equation}
%where $C$ is the Cartier operator.
%\end{prop}
\begin{proof}
Choose the affine open part $x_0 = 1$ of the moduli for Calabi-Yau n-folds. Based on the computational result of Cartier operators of Miller in \cite[Corollary 1]{Miller72}, we get
\begin{equation}
C(\omega) = \frac{\psi(f^{p-1}x_1x_2\cdots x_{n+1})}{x_1x_2\cdots x_{n+1}} \omega,
\end{equation}
where $\psi$ the $p^{-1}$-linear operator defined over $\mathbb{Q}[t_1, t_{n+2}, \frac{1}{(n+2)(t_1^{n+2}-t_{n+2})}][x_1,x_2,\cdots,x_{n+1}]$ given by
\[
\psi(x_1^{i_1}x_2^{i_2}\cdots x_{n+1}^{i_{n+1}}) = \begin{cases}
x_1^{i_1/p}x_2^{i_2/p}\cdots x_{n+1}^{i_{n+1}/p}, & \mathrm{if} \ p|i_j, j=1,2,\cdots, i_{n+1}; \\
0, & \mathrm{otherwise}.
\end{cases}
\]
Note that the coefficient $(x_1x_2\cdots x_{n+1})^{p-1}$ in $f^{p-1}$ is
\begin{equation}
    \begin{split}
  &\sum_{k=0}^{\lfloor \frac{p-1}{n+2} \rfloor}\frac{(p-1)!}{(k!)^{n+2}(p-1-(n+2)k)!} (-1)^{(n+2)k}t_{n+2}^k ((n+2)t_1)^{p-1-(n+2)k} \\
  =& \sum_{k=0}^{\lfloor \frac{p-1}{n+2} \rfloor}\frac{((n+2)k)!}{(k!)^{n+2}}
\binom{p-1}{(n+2)k} (-1)^{(n+2)k} t_{n+2}^k ((n+2)t_1)^{p-1-(n+2)k}.      
    \end{split}
\end{equation}

Together with the identity 
$$
\binom{p-1}{(n+2)k}\equiv_p (-1)^{(n+2)k},
$$
we get the final result.
\end{proof}



















\begin{comment}
\begin{coro}
If we let $X =5t_0$ and $Y = t_4$, then
\begin{equation}
A(X, Y) = \sum_{i=0}^{\lfloor \frac{p-1}{5} \rfloor} \frac{(p-1)!}{(i!)^5(p-1-5i)!} (-Y)^i X^{p-1-5i}. 
\end{equation}
\end{coro}
For any given prime $p$, one could use the above singular code to compute $A(x,y)$. For example, when $p=13$, the codes are as follows:
\begin{tiny}
\begin{verbatim}
int p = 13;  // CHANGE THIS TO ANY PRIME YOU WANT
execute("ring R = " + string(p) + ", (X,Y), dp;");
int max_i = (p-1) div 5;
poly f = 0;
int i, j;
number c, d1, d2;
i = 0;
while (i <= max_i)
{c = -1;  // Wilson's theorem
    // Compute i!^5
    d1 = 1;
    j = 1;
    while (j <= i)
    {d1 = d1 * j;j = j + 1;}
    d1 = d1^5;
    // Compute (p-1-5i)!
    d2 = 1;
    j = 1;
    while (j <= (p-1-5*i))
    {d2 = d2 * j; j = j + 1;}
    c = c / (d1 * d2);
    // Apply sign from (-Y)^i
    if (i % 2 == 1)
    {c = -c;}
    if (c != 0)
    {f = f + c * X^(p-1-5*i) * Y^i;}
    i = i + 1;}
"Polynomial for p = " + string(p) + ":";
f; 
\end{verbatim}    
\end{tiny}

Our conventions in this section will be such that the conifold is at $z=1$. 

\textbf{To do:} Compute the transition matrices for the periods and find the behavior of the holomorphic period at the Fermat point. $t_4$ should be some combination of $z$ and the holomorphic period that is regular at the Fermat point.
\end{comment}
The logarithmic period is given by 
$$
(n+2)(2\pi i )^{-n}\int_{\delta_{1,z}}\alpha_z=F(z)
\ln(\frac{z}{(n+2)^{n+2}})+G(z),\ \ 
$$
$$
G(z):= (n+2)\sum_{k=1}^{\infty}  \frac{((n+2)k)!}{(k!)^{n+2}} (\sum_{j=k+1}^{(n+2)k}\frac{1}{j}) (\frac{z}{(n+2)^{n+2}})^k.
$$%\marginpar{\tiny {\color{red} Jin} The exact expressions here, see the article \cite{GMPR}.  }
The period expression of $t_1$ and $t_{n+1}$ up to some $2\pi i$ factors
are 
$$
t_1=F(z),\ t_{n+2}=zF(z)^{n+2}. 
$$
After inserting mirror map $z(q)$, we get the $q$-expansion of these quantities. For example, when $n=1$, we get the reversion formula:
\[
t_1(q) = \ _2F_1(\frac{1}{3},\frac{2}{3},1;z(q)) = \theta_3(q)\theta_3(q^3)+\theta_2(q)\theta_2(q^3)
\]
and
\[ t_3(q) = z(q) _2F_1(\frac{1}{3},\frac{2}{3},1;z(q)) = 27\frac{\eta^9(q^3)}{\eta^3(q)},\]
whose $q$-expansions are:
$$
t_1(q)= 1+6q+6q^3+6q^4+O\left(q^6\right),\ \   t_3(q)= 27q+81q^2+243q^3+351q^4+729q^5+O\left(q^6\right).\ \ 
$$
See \cite[Section 8]{GMCD-DF} and \cite[Section 5.1]{Younes2019}. These formula has been shown in \cite[Corollary 2.4]{BBG1994} via the cubic theta functions.

For $n=2$, the mirror map is given by
\[
z(q) = q-104q^2+6444q^3-311744q^4+13018830q^5 + O(q^6),
\]
which can be found in \cite[(7.3.7)]{Gannon2023} and \cite[A286329]{oeis}.  Then we have:
$$
t_1(q)=\ _3F_2(\frac{1}{4},\frac{2}{4}, \frac{3}{4};1,1| z(q)) =  \theta_{D_4}^4(q) = 1+24q+24q^2+96q^3+24q^4+144q^5+O(q^6),   
$$
where $\theta_{D_4}(q)$ is the theta series of $D_4$-lattice and 
$$
t_4(q) = z(q) _3F_2(\frac{1}{4},\frac{2}{4}, \frac{3}{4};1,1| z(q))^4 = c \eta^8(q) \eta^8(q^2) = c(q-8q^2+12q^3+64q^4-210q^5+O(q^6)),
$$
where $\eta(q)$ is the eta function and $c$ is a constant. %\marginpar{\tiny Mohamod, could you use the mirror map $z(q)$ and hypergeometric functions to fix these constants?} 
For the q-expansion of $t_1(q)$ (resp. $t_4(q)$), see \cite[A004011]{oeis} (resp. \cite[A002288]{oeis}). These results are coincide with the q-expansion computations in \cite{GMCD-DF}. 
\

For $n=3$ we get
\begin{align}
    \begin{split}
        t_1(q)&=1+120 q+21000 q^2+14115000 q^3+13414125000 q^4+15234972675120
   q^5+O\left(q^6\right)\\
   t_5(q) &= q-170 q^2-41475 q^3-32183000 q^4-32678171250 q^5+O\left(q^6\right)\\
    \end{split}
\end{align}
which appear in \cite{GMCD-MQCY3}. 

For $n=4$ we get 
$$
t_1=\frac{1}{6}+20q+82620q^2+O(q^3),\   t_6=\frac{1}{46656}q-\frac{1}{24}q^2 + O(q^2),\ \ 
$$
see \cite[Section 8.3]{GMCD-DF}.%\marginpar{\tiny {\color{red} Mohamed}  }

For the purpose of studying supersingular Calabi-Yau threefolds, we have also listed some Hasse-Witt invariants of the mirror quintic in Table~\ref{table: A_p(x,y) of mirror quintic at small primes} as product of irreducible polynomials over $\mathbb{F}_p[x,y]$, where $x:=t_1,\ y:=t_{n+2}$. 
\begin{table}
\begin{equation*}
\begin{array}{|c|l|}
\hline
 p & A_p(x,y) \\
 \hline
 2 & x \\[2pt]
 3 & x^2 \\[2pt]
 5 & x^4 \\[2pt]
 7 & x \left(x^5+y\right) \\[2pt]
 11 & x^{10}+10 x^5 y+y^2 \\[2pt]
 13 & x^2 \left(x^{10}+3 x^5 y+y^2\right) \\[2pt]
 17 & x \left(x^5+16 y\right) \left(x^{10}+2 x^5 y+12 y^2\right) \\[2pt]
 19 & x^3 \left(x^{15}+6 x^{10} y+8 x^5 y^2+7 y^3\right) \\[2pt]
 23 & x^2 \left(x^5+14 y\right) \left(x^5+21 y\right) \left(x^{10}+16 x^5 y+7 y^2\right) \\[2pt]
 29 & x^3 \left(x^{10}+15 y^2\right) \left(x^{15}+4 x^{10} y+24 x^5 y^2+14 y^3\right) \\[2pt]
 31 & \left(x^5+4 y\right) \left(x^5+17 y\right) \left(x^{20}+6 x^{15} y+25 x^{10} y^2+3 x^5 y^3+y^4\right) \\[2pt]
 37 & x \left(x^5+9 y\right) \left(x^5+23 y\right) \left(x^{25}+14 x^{20} y+6 x^{15} y^2+21 x^{10} y^3+17 x^5 y^4+3 y^5\right) \\[2pt]
 41 & \left(x^5+5 y\right) \left(x^{10}+26 x^5 y+27 y^2\right) \left(x^{25}+7 x^{20} y+30 x^{15} y^2+31 x^{10} y^3+12 x^5 y^4+31 y^5\right) \\[2pt]
 43 & x^2 \left(x^{40}+34 x^{35} y+9 x^{30} y^2+31 x^{25} y^3+20 x^{20} y^4+2 x^{15} y^5+41 x^{10} y^6+11 x^5 y^7+29 y^8\right) \\[2pt]
 47 & x \left(x^5+21 y\right) \left(x^{10}+23 x^5 y+17 y^2\right) \left(x^{10}+44 x^5 y+23 y^2\right) \left(x^{20}+32 x^{15} y+25 x^{10} y^2+15 x^5 y^3+3 y^4\right) \\[2pt]
 \hline
 % 53 & x^2 \left(x^5+32 y\right) \left(x^5+50 y\right) \left(x^{40}+38 x^{35} y+34 x^{30} y^2+19 x^{25} y^3+28 x^{20} y^4+2 x^{15} y^5+34 x^{10} y^6+35 x^5 y^7+41 y^8\right) \\
 % 59 & x^3 \left(x^5+49 y\right) \left(x^{10}+3 x^5 y+38 y^2\right) \left(x^{40}+9 x^{35} y+57 x^{30} y^2+4 x^{25} y^3+33 x^{20} y^4+4 x^{15} y^5+44 x^{10} y^6+46 x^5 y^7+40
 %   y^8\right) \\
 % 61 & \left(x^5+59 y\right) \left(x^{15}+6 x^{10} y+41 x^5 y^2+11 y^3\right) \left(x^{40}+55 x^{35} y+57 x^{30} y^2+45 x^{25} y^3+49 x^{20} y^4+13 x^{15} y^5+19 x^{10} y^6+10 x^5
 %   y^7+20 y^8\right) \\
 % 67 & x \left(x^5+6 y\right) \left(x^5+23 y\right) \left(x^{55}+24 x^{50} y+6 x^{45} y^2+8 x^{40} y^3+7 x^{35} y^4+8 x^{30} y^5+47 x^{25} y^6+29 x^{20} y^7+49 x^{15} y^8+48 x^{10}
 %   y^9+15 x^5 y^{10}+65 y^{11}\right) \\
 % 71 & \left(x^{10}+67 x^5 y+34 y^2\right) \left(x^{60}+53 x^{55} y+49 x^{50} y^2+54 x^{45} y^3+9 x^{40} y^4+56 x^{35} y^5+18 x^{30} y^6+13 x^{25} y^7+46 x^{20} y^8+23 x^{15} y^9+18
 %   x^{10} y^{10}+6 x^5 y^{11}+20 y^{12}\right) \\
\end{array}
\end{equation*}
\caption{ $A_p(x,y)$ of mirror quintic at small primes $p$.}
\label{table: A_p(x,y) of mirror quintic at small primes}
\end{table}


In a similar way we compute the Hasse-Witt invariant of the rest three families of CY threefolds in \cref{28112025bri}. %\marginpar{\tiny \color{red} Jin: Write the computation of four cases here.}


\begin{prop} 
\begin{enumerate}
    \item 
For the hypersurface in the weighted projective space $\mathbb{P}^{5,2,1,1,1}$ defined by
\[
-x_0^2 - x_1^5 - t_{10}x_2^{10} - x_3^{10} - x_4^{10} +t_1 x_0x_1x_2x_3x_4 = 0,
\]
we have:
\begin{equation}
C(\alpha) = \left( \sum_{k=0}^{\lfloor \frac{p-1}{10} \rfloor} \frac{(10k)!}{(k!)^3(2k)!(5k)!}t_{10}^k t_1^{p-1-10k} \right)^{\frac{1}{p}} \alpha,
\end{equation}
where $p$ is a prime number such that $p\not|10$.
\item
For the hypersurface in the weighted projective space $\mathbb{P}^{4,1,1,1,1}$ defined by
\[
-x_0^2 -t_8 x_1^8 - x_2^{8} - x_3^{8} - x_4^{8} +t_1 x_0x_1x_2x_3x_4 = 0,
\]
we have:
\begin{equation}
C(\alpha) = \left( \sum_{k=0}^{\lfloor \frac{p-1}{8} \rfloor} \frac{(8k)!}{(k!)^4(4k)!}t_{8}^k t_1^{p-1-10k} \right)^{\frac{1}{p}} \alpha,
\end{equation}
where $p$ is a prime number such that $p\not| 2$.
\item 
For the hypersurface in the weighted projective space $\mathbb{P}^{1,1,1,2,1}  $ defined by
\[
-t_6x_0^6 - x_1^6 - x_2^{6} - x_3^{3} - x_4^{6} +t_1 x_0x_1x_2x_3x_4 = 0,
\]
we have:
\begin{equation}
C(\alpha) = \left( \sum_{k=0}^{\lfloor \frac{p-1}{6} \rfloor} \frac{(6k)!}{(k!)^4(2k)!}t_{6}^k t_1^{p-1-6k} \right)^{\frac{1}{p}} \alpha,
\end{equation}
where $p$ is a prime number such that $p\not|6$.
\end{enumerate}
\end{prop}


\begin{proof}
The proof is similar to the proof in Proposition 1. Let's take the first one as an example. Choose the affine open part $x_0 = 1$ of the first family for Calabi-Yau 3-folds. Based on the formula
\begin{equation}
C(\omega) = \frac{\psi(f^{p-1}x_1x_2x_3x_{4})}{x_1x_2x_3x_{4}} \omega,
\end{equation}
we only need to compute the coefficient $(x_1x_2x_3x_4)^{p-1}$ in $f^{p-1}$, which is
\[
\sum_{k=0}^{\lfloor \frac{p-1}{10} \rfloor}\frac{(p-1)!}{(k!)^{3}(2k)!(5k)!(p-1-(10k)!)} (-1)^{10k}t_{10}^k t_1^{p-1-10k}= \sum_{k=0}^{\lfloor \frac{p-1}{10} \rfloor}\frac{(10k)!}{(k!)^{3}(2k)!(5k)!} t_{10}^k t_1^{p-1-10k}.
\]
Hence we have:
\[
C(\alpha) = \left( \sum_{k=0}^{\lfloor \frac{p-1}{10} \rfloor} \frac{(10k)!}{(k!)^3(2k)!(5k)!}t_{10}^k t_1^{p-1-10k} \right)^{\frac{1}{p}} \alpha
\]
in this case.
\end{proof}
\begin{rem}
We also have the following identities:
\begin{eqnarray*}
\frac{(\frac{1}{5})_k (\frac{4}{5})_k (\frac{2}{5})_k (\frac{3}{5})_k }{k!^4} &=& \frac{1}{5^{5k}} \frac{(5k)!}{(k!)^5} \\
\frac{(\frac{1}{10})_k (\frac{9}{10})_k (\frac{3}{10})_k (\frac{7}{10})_k }{k!^4} &=& \frac{1}{10^{3k}5^{2k}2^{5k}} \frac{(10k)!}{(k!)^3 (2k)! (5k)!} \\
\frac{(\frac{1}{8})_k (\frac{7}{8})_k (\frac{3}{8})_k (\frac{5}{8})_k }{k!^4} &=& \frac{1}{8^{4k}2^{4k}} \frac{(8k)!}{(k!)^4 (4k)!} \\
\frac{(\frac{1}{6})_k (\frac{5}{6})_k (\frac{2}{6})_k (\frac{4}{6})_k }{k!^4} &=&\frac{1}{6^{4k}3^{2k}} \frac{(6k)!}{(k!)^4 (2k)!}.
\end{eqnarray*}
\end{rem}
    




\section{Calabi-Yau Operators}

The function $A_p(z)$, as defined in \cref{17112025ymsc}, can be computed directly from the periods. It follows that specific cases of the conjecture can be tested without any reference to an underlying family of varieties. We test this conjecture for a list of 545 Calabi-Yau operators and find that it fails sometimes. This list of Calabi-Yau opearators consists mainly of operators from the AESZ list \cite{alenstzu} and can be found as an attached file, along with a SageMath Jupyter notebook.

A Calabi-Yau operator is a fourth order linear differential operator in a variable $z$, which satisfies a number of properties that make it a candidate Picard-Fuchs equation for a family of Calabi-Yau threefolds \cite{alenstzu,almzud,vanStratenCYOperators}. These properties include the requirement that $z=0$ is a point of maximal unipotent monodromy (MUM), the mirror map $z(q)$ defined by this MUM point has integer coefficients when expanded in $q$, the genus $0$ instanton numbers are integers up to some overall multiple, etc.

There is no guarantee that all of the Calabi-Yau operators in the AESZ list \cite{alenstzu} are really the Picard-Fuchs operator for a family of Calabi-Yau threefolds with $h^{2,1}=1$, so these examples do not immediately invalidate Conjecture~\ref{17112025ymsc}. It is tempting to speculate that, for a given Calabi-Yau operator, the failure of Conjecture~\ref{17112025ymsc} rules out the existence of an associated smooth Calabi-Yau threefold with $h^{2,1}=1$.
%\marginpar{\tiny Is there an example of a Calabi-Yau operator, which has been shown that it doesn't come from a smooth family of Calabi-Yau threefolds with $h^{2,1}=1$?}
%\marginpar{\tiny {\color{red}  I'm not sure what's been proven, but there are operators for which the standard recipe for finding a basis of periods with $Sp(4,\mathbb{Z})$ monodromy fails (see \cite{vanvan}). This strongly suggests that there is no smooth family of mirror CY threefolds with $h^{2,1}=1$. I can do some monodromy calculations when I'm back in Beijing.}}

% For each opearotr, we compute the holomorphic period 
% \begin{equation}
%     {\rm HP}(X_z,\alpha_z)=\sum_{n=0}^{\infty}c_nz^n
% \end{equation}
% around $z=0$, which has been normalised so that $c_0=1$. The mirror map $z(q)$ to a high order in $q$ by inverting the quotient of the logarithmic period and the holomorphic period.

For each operator, we compute the holomorphic period
\begin{equation*}
    \varpi_0(z) = \sum_{n=0}^{\infty}c_nz^n
\end{equation*}
%\marginpar{\tiny $\varpi_0$ is probably better than ${\rm HP}(X_z,\alpha_z)$ for the holomorphic period in this section, because we don't know that a $X_z$ exists for a given operator. }
around $z=0$ and normalise it so that $c_0=1$. 


% We also compute the mirror map $z(q)$ to a high order in $q$ by inverting the quotient of the logarithmic period and the holomorphic period. 

From a list of 545 Calabi-Yau operators, we find that 460 operators satisfy
% \begin{equation}
%     A_p(z(q)) =\left(\sum_{m=0}^{\infty}c_m z(q)^m\right)^{p-1}~\sum_{n=0}^{p-1}c_nz(q)^n    \equiv_p1 +O(q^{200})
% \end{equation}
\begin{equation*}
    A_p(z) =\varpi_0(z)^{p-1}~\sum_{n=0}^{p-1}c_nz^n    \equiv_p1 +O(z^{600})
\end{equation*}
for the first $100$ primes. The truncation at $O(z^{600})$ is an arbitrarily chosen cutoff.

The remaining 85 operators do not satisfy this condition. %\marginpar{\tiny List this 85} 
However, we still find that
\begin{equation*}
    A_p(z) \equiv_p1 +O(z^{p})
\end{equation*}
for all of the operators that we check and the first $100$ primes.



When Conjecture~\ref{17112025ymsc} fails, it often fails in interesting ways. For example, consider the operator 
\begin{align*}
\label{eq: AESZ 1hat after change of variables}
    \begin{split}
        \mathcal{L}=&\theta^4-2\cdot 5(10000\theta^4 + 12500\theta^3 + 9500\theta^2 + 3250\theta + 399)\\
        &+2^2\cdot5^8z^2(2400\theta^4 + 6000\theta^3 + 6290\theta^2 + 2800\theta + 399)\\
        &-2^5\cdot 5^{14}z^3(4\theta + 3)(80\theta^3 + 240\theta^2 + 221\theta + 42)\\
        &+2^4\cdot5^{20}z^4(4\theta + 1)(4\theta + 3)(4\theta + 7)(4\theta + 9)
    \end{split}
\end{align*}
where $\theta=z\frac{d}{dz}$. This operator can be obtained from operator $\widehat{1}$ in \cite{alenstzu} by making the change of variables $z\mapsto 2z$. The change of variables is needed to ensure that the holomorphic period has integer coefficients.

The singularities of this operator are summarised by the Riemann $\mathcal{P}$-symbol
\begin{equation*}
    \mathcal{P}\left\{\begin{array}{c c c} 0 & \frac{1}{2^3\cdot 5^5} & \infty\\[1pt]
    \hline
    0 & \frac{1}{20} & \frac{1}{4}\\
    0 & \frac{3}{20} & \frac{3}{4}\\
    0 & \frac{7}{20} & \frac{7}{4}\\
    0 & \frac{9}{20} & \frac{9}{4}\\
    \end{array}\right\},
\end{equation*}
which is somewhat exotic because it does not have a conifold point (a singularity with indices $(0,1,1,2)$).

The holomorphic period for this operator is given by
\begin{align*}
    \begin{split}
        \varpi_0(z)=&1+3990 z+49934850 z^2+806586406500 z^3+14635392749853750 z^4+O\left(z^{5}\right)
    \end{split}
\end{align*}
% \begin{align}
%     \begin{split}
%         \varpi_0(z)=&1+3990 z+49934850 z^2+806586406500 z^3+14635392749853750 z^4+O\left(z^{5}\right)\\
%    z(q)=&q-16540 q^2+232386100 q^3-3003408484000 q^4+36919670024260000 q^5+O\left(q^{6}\right),
%     \end{split}
% \end{align}
which we use to compute $A_p(z)$ for many primes and to a high order in $z$.

After some experimentation, we are led to the following conjecture, which we verify for the first $120$ primes and to $O(z^{10000})$.

\begin{conj}
\label{01022026aquamatica}
If $p$ is an inert rational prime in $\mathbb{Q}\left(\sqrt{-5}\right)$, then
\begin{equation*}
    A_p(z)\equiv_p \left(\sqrt{1-2^3\cdot5^5 z}\right)^p
\end{equation*}
where the right hand side should be understood as a power series in $z$.
\end{conj}

Conjecture~\ref{01022026aquamatica} is a result of computer experiments. In order to test Conjecture~\ref{01122025zutopia}, we  computed $A_p(z)$ for many primes and to a high order in $z$. A few examples at small primes are listed in Table~\ref{table: A_p(z) at small primes for pullback of operator 1hat}.  It quickly becomes clear that, for half of the primes, $A_p(z)$ contains only $p^{th}$ powers of $z$. In other wrds,
\begin{equation*}
    A_p(z)\equiv_p f_p(z)^p.
\end{equation*}
for some $f_p(z)=\mathbb{F}_p\llbracket z\rrbracket$.
With enough primes, we are able to conjecture that $A_p(z)$ is a $p^{th}$ power whenever $p$ is inert in $\mathbb{Q}\left(\sqrt{-5}\right)$.

At this point, one naturally wonders if the $f_p(z)$ are the mod $p$ reduction of some $f(z)\in\mathbb{Z}\llbracket z\rrbracket$. It is straightforward to search for a candidate $f$ by applying the Chinese remainder theorem term-wise to $f_p$ for many primes. This defines an element of $\left(\mathbb{Z}/N\mathbb{Z}\right)\llbracket z\rrbracket$ for some large $N\in\mathbb{N}$. We then apply the rational reconstruction algorithm term-wise to this power series to find a power series with integer coefficients, which stabilises as we consider more primes and increase $N$. We recognize the resulting power series as the power series expansion of $\sqrt{1-2^3\cdot5^5 z}$. The computer data used in this section can be found in \href{https://w3.impa.br/~hossein/WikiHossein/files/Singular%20Codes/25_02_2026_HW_For_CY/}{third author's webpage.}
%\footnote{https://w3.impa.br/$\sim$hossein/WikiHossein/files/Singular codes/25_02_2026_HW_For_CY/}

Observations similar to Conjecture~\ref{01022026aquamatica} have appeared in \cite{Ducker:2025wfl}, where the authors study a number of operators like the operator $\mathcal{L}$ in equation~\eqref{01022026aquamatica}. These operators are obtained as pullbacks of fifth-order differential operators \cite{almzud,2006math.....12215A}. 

The authors of \cite{Ducker:2025wfl} note that, in specific examples, the fourth order operator cannot have a basis of solutions with rational monodromy and that this is remedied by passing to a double cover of the $z$-plane with quadratic branch points introduced at the two non-zero singularities. In our example, this would make the function $\sqrt{1-2^3\cdot5^5 z}$ single valued. They also note that, without passing to the double cover, the deformation method for computing local zeta functions from Picard-Fuchs equations fails \cite{Candelas:2021tqt}. These observations support the idea that the failure of Conjecture~\ref{17112025ymsc} rules out the existence of an associated smooth Calabi-Yau threefold with $h^{2,1}=1$.

% An observation similar to Conjecture~\ref{01022026aquamatica} has been made in \cite{Ducker:2025wfl} by studying a fourth order differential operator similar to the operator $\mathcal{L}$ in \eqref{eq: AESZ 1hat after change of variables}. Both operators have a MUM point at $0$ and two additional singularities with finite monodromy. The authors of \cite{Ducker:2025wfl} observe that, while the initial operator cannot have a basis of solutions with rational monodromy, one finds another operator that does after a moving to a double cover of the $z$-plane. 

Finally, we note that a single Calabi-Yau operator can have multiple MUM points. The standard example is the Picard-Fuchs operator described by Rodland in \cite{1998math......1092R}. This operator appears as operator 27 in \cite{almzud} and is given by
\begin{align*}
    \mathcal{L}_{Rodland}=&3^{2} \theta^4-3 z\left(173\theta^4+340\theta^3+272\theta^2+102\theta+15\right)\\
    &-2 z^{2}\left(1129\theta^4+5032\theta^3+7597\theta^2+4773\theta+1083\right)\\
    &+2 z^{3}\left(843\theta^4+2628\theta^3+2353\theta^2+675\theta+6\right)\\
    &-z  ^{4}\left(295\theta^4+608\theta^3+478\theta^2+174\theta+26\right)\\
    &+z^{5}\left(\theta+1\right)^4
\end{align*}
It has the Riemann symbol
\begin{equation*}
    \mathcal{P}\left\{\begin{array}{c c c c} -3 & 0 & \omega_i & \infty\\[1pt]
    \hline
    0 & 0 & 0 & 1\\
    1 & 0 & 1 & 1\\
    3 & 0 & 1 & 1\\
    4 & 0 & 2 & 1\\
    \end{array}\right\},
\end{equation*}
\noindent where $\omega_i$ is any one of the solutions of $\omega_i^3-289\omega_i^2-57\omega_i+1=0$. There are MUM points at $z=0$ and $z=\infty$ with holomorphic solutions 
\begin{align*}
    \begin{split}
        \varpi_0^{(0)}(z)&=1+5 z+109 z^2+3317 z^3+121501 z^4+4954505
   z^5+216867925 z^6+O\left(z^{7}\right)\\
   \varpi_0^{(\infty)}(z)&=1+ \frac{17}{z}+ \frac{1549}{z^2}+ \frac{215585}{z^3}+ \frac{36505501}{z^4}+
   \frac{6921832517}{z^5}+ \frac{1412721479989}{z^6}+O\left(z^{7}\right)
    \end{split}
\end{align*}
We find that the Hasse-Witt invariants computed at each of these MUM points are related by $z\rightarrow\frac{1}{z}$. For example, at $p=7$, we find
\begin{align*}
    A_7^{(0)}(z)&=z^6+4 z^5+2 z^4+2 z^2+4 z+1\\
    A_7^{(\infty)}(z)&=\frac{1}{z^6}+\frac{4}{z^5}+\frac{2}{z^4}+\frac{2}{z^2}+\frac{4} {z}+1~.
\end{align*}





\begin{table}
\begin{equation*}
\begin{array}{|c|l|}
\hline
 p & A_p(z) \\[2pt]
 \hline
2 & 1 + O(z^{100}) \\[4pt]
3 & 1 + O(z^{100}) \\[4pt]
5 & 1 + O(z^{100}) \\[4pt]
7 & 1 + O(z^{100}) \\[4pt]
11 & 1 + 7 z^{11} + 3 z^{22} +  z^{33} + 5 z^{44} + 6 z^{55} + 3 z^{66} + O(z^{100}) \\[4pt]
13 & 1 + 6 z^{13} + 8 z^{26} + 4 z^{39} + 9 z^{52} + 5 z^{65} + 7 z^{78} + 12 z^{91} + O(z^{100}) \\[4pt]
17 & 1 + 12 z^{17} + 13 z^{34} + 14 z^{51} + 11 z^{68} + 9 z^{85} + O(z^{100}) \\[4pt]
19 & 1 + 2 z^{19} + 17 z^{38} + 4 z^{57} + 9 z^{76} + 9 z^{95} + O(z^{100}) \\[4pt]
23 & 1 + 22 z^{23} + 12 z^{24} + 15 z^{25} + 7 z^{26} +  z^{27} + 8 z^{29} + 20 z^{30} + 15 z^{31} + 12 z^{32} + 5 z^{33} + 22 z^{34} + 19 z^{35}\\
& ~+ 15 z^{36} + 9 z^{37} + 21 z^{38} + 15 z^{39} +  z^{40} + 22 z^{41} + 3 z^{42} + 3 z^{43} +  z^{44} + 8 z^{45} + 4 z^{46} + 12 z^{47} + 19 z^{48}\\
& ~+  z^{49} + 6 z^{50} + 11 z^{51} + 18 z^{52} +  z^{53} + 11 z^{54} + 18 z^{55} + 5 z^{56} +  z^{57} + 4 z^{58} + 8 z^{59} + 9 z^{60} + 2 z^{61} + 19 z^{62}\\
&~+ 15 z^{63} + 19 z^{64} + 2 z^{65} + 16 z^{66} + 4 z^{67} + 18 z^{68} + 7 z^{69} + 4 z^{70} +  z^{71} + 21 z^{72} + 11 z^{73} + 6 z^{74} + 13 z^{75}\\
& ~+ 18 z^{76} + 16 z^{77} + 21 z^{78} + 21 z^{79} + 8 z^{80} + 3 z^{81} + 6 z^{82} + 10 z^{83} + 15 z^{84} + 22 z^{85} + 4 z^{86} + 17 z^{87} +  z^{88}\\
& ~+ 10 z^{89} + 22 z^{90} + 18 z^{91} + 9 z^{92} + 19 z^{93} + 18 z^{94} + 18 z^{96} + 21 z^{97} +  z^{98} + 8 z^{99} + O(z^{100}) \\[4pt]
29 & 1 + 13 z^{29} + 8 z^{30} + 5 z^{31} + 21 z^{32} + 24 z^{33} + 15 z^{34} + 17 z^{35} + 2 z^{36} + 28 z^{37} + 8 z^{38} + 24 z^{39} + 12 z^{40}\\
&~+ 18 z^{41} + 20 z^{43} + 16 z^{44} + 15 z^{45} + 26 z^{46} + 22 z^{47} + 4 z^{48} + 22 z^{49} + 15 z^{50} + 15 z^{51} + 26 z^{52} + 28 z^{53}\\
&~+ 3 z^{54} + 8 z^{55} + 26 z^{56} + 14 z^{57} + 9 z^{58} + 5 z^{59} +  z^{60} + 13 z^{61} + 14 z^{62} + 9 z^{63} + 2 z^{64} + 10 z^{65} + 21 z^{66}\\
&~+ 24 z^{67} + 10 z^{68} + 10 z^{70} + 26 z^{71} + 17 z^{72} + 14 z^{73} + 2 z^{74} +  z^{75} + 16 z^{76} + 10 z^{77} + 8 z^{78} + 22 z^{79} + 7 z^{80}\\
&~+ 11 z^{81} + 9 z^{82} + 8 z^{83} + 2 z^{84} + 8 z^{85} + 21 z^{86} + 4 z^{87} + 27 z^{88} + 18 z^{89} + 4 z^{90} + 14 z^{91} + 21 z^{92} + 19 z^{93}\\
&~ + 28 z^{94} + 12 z^{95} + 25 z^{96} + 6 z^{97} + 14 z^{98} + 9 z^{99} + O(z^{100}) \\[4pt]
31 & 1 + 24 z^{31} + 22 z^{62} + 30 z^{93} + O(z^{100}) \\[4pt]
37 & 1 + 6 z^{37} + 19 z^{74} + O(z^{100}) \\[4pt]
41 & 1 + 34 z^{41} + 40 z^{42} + 19 z^{43} + 15 z^{45} + 28 z^{46} + 28 z^{47} + 33 z^{48} + 36 z^{49} + 5 z^{50} + 40 z^{51} + 19 z^{52} + 28 z^{53}\\
& ~+ 38 z^{54} + 33 z^{55} + 13 z^{56} + 22 z^{57} + 25 z^{58} + 21 z^{59} + 38 z^{60} + 11 z^{61} + 39 z^{62} + 31 z^{63} +  z^{64} + 19 z^{65}\\
&~ + 31 z^{66} + 10 z^{67} + 8 z^{68} + 4 z^{69} + 6 z^{70} + 35 z^{71} + 32 z^{72} + 9 z^{73} + 39 z^{74} + 38 z^{75} +  z^{76} + 12 z^{77} + 4 z^{78}\\
&~+ 37 z^{79} + 30 z^{80} + 9 z^{81} + 4 z^{82} + 38 z^{83} + 29 z^{84} + 16 z^{85} + 29 z^{86} + 35 z^{87} + 5 z^{88} + 12 z^{89} + 25 z^{90} + 35 z^{91}\\
&~ + 24 z^{92} + 27 z^{93} + 6 z^{94} + 27 z^{95} + 33 z^{96} + 8 z^{97} + 10 z^{98} + 25 z^{99} + O(z^{100}) \\[4pt]
43 & 1 + 40 z^{43} + 35 z^{44} + 11 z^{45} + 22 z^{46} + 34 z^{47} + 20 z^{48} + 6 z^{49} + 40 z^{50} + 29 z^{51} + 14 z^{52} +  z^{53} + 28 z^{54}\\
&~+ 30 z^{55} + 14 z^{56} + 16 z^{57} + 5 z^{58} +  z^{59} + 12 z^{60} + 40 z^{61} + 7 z^{62} + 7 z^{63} + 25 z^{64} + 18 z^{65} + 34 z^{66} + 9 z^{67}\\
&~+ 27 z^{68} + 28 z^{69} + 19 z^{70} + 6 z^{71} + 14 z^{72} + 29 z^{73} + 29 z^{74} + 21 z^{75} + 23 z^{76} + 8 z^{77} + 12 z^{78} + 39 z^{79} + 35 z^{80}\\
&~+ 33 z^{81} + 26 z^{82} + 7 z^{83} + 27 z^{84} + 30 z^{85} + 38 z^{86} + 42 z^{87} + 10 z^{88} + 6 z^{89} + 42 z^{90} + 34 z^{91}+ 36 z^{92} + 4 z^{93}\\
&~+ 20 z^{94} + 4 z^{95} + 25 z^{96} + 26 z^{97} + 10 z^{98} + 20 z^{99} + O(z^{100}) \\[4pt]
47 & 1 + 17 z^{47} + 22 z^{48} + 37 z^{49} + 46 z^{50} + 5 z^{51} + 10 z^{52} + 32 z^{53} + 2 z^{54} + 38 z^{55} + 25 z^{56} + 32 z^{57} + 28 z^{58}\\
&~+ 27 z^{59} + 22 z^{60} + 25 z^{61} + 40 z^{62} + 37 z^{63} + 31 z^{64} + 15 z^{65} + 13 z^{66} + 45 z^{67} + 11 z^{68} + 26 z^{69} + 18 z^{70}\\
&~+ 20 z^{71} + 10 z^{72} + 5 z^{73} + 3 z^{74} + 26 z^{75} + 3 z^{76} + 5 z^{77} + 40 z^{78} + 37 z^{79} + 45 z^{80} + 4 z^{81} + 22 z^{82} +  z^{83}\\
&~+ 39 z^{84} + 7 z^{85} + 3 z^{86} + 4 z^{87} + 3 z^{88} + 21 z^{89} + 33 z^{90} + 27 z^{91} + 20 z^{92} + 30 z^{94} + 15 z^{95} + 10 z^{96} + 24 z^{97}\\
&~+ 13 z^{98} + 43 z^{99} + O(z^{100}) \\[4pt]
53 & 1 + 8 z^{53} + O(z^{100}) \\[4pt]
59 & 1 + 8 z^{59} + O(z^{100}) \\[4pt]
% 61 & 1 + 45 z^{61} + 60 z^{62} + 56 z^{63} + 6 z^{64} + 60 z^{65} + 34 z^{66} + 29 z^{67} + 52 z^{68} + 10 z^{69} + 34 z^{70} + 45 z^{71} +  z^{72} + 26 z^{73} + 50 z^{74} + 49 z^{75} + 22 z^{76} +  z^{77} + 13 z^{78} + 40 z^{79} + 13 z^{80} + 31 z^{81} + 37 z^{82} + 19 z^{83} + 48 z^{84} + 10 z^{85} + 48 z^{86} + 55 z^{87} + 24 z^{88} + 18 z^{89} + 24 z^{90} + 52 z^{91} + 37 z^{92} + 17 z^{93} + 10 z^{94} + 5 z^{95} + 39 z^{96} + 23 z^{97} + 16 z^{98} + 14 z^{99} + O(z^{100}) \\[4pt]
% 67 & 1 + 37 z^{67} + 28 z^{68} + 7 z^{69} + 57 z^{70} + 50 z^{71} + 8 z^{72} + 43 z^{74} + 19 z^{75} + 7 z^{76} + 66 z^{77} + 37 z^{78} + 31 z^{79} + 29 z^{80} + 30 z^{81} + 66 z^{82} + 27 z^{83} + 15 z^{84} + 44 z^{85} + 45 z^{86} + 12 z^{87} + 49 z^{88} + 50 z^{89} + 50 z^{90} + 63 z^{91} + 42 z^{92} + 43 z^{93} + 35 z^{94} + 5 z^{95} + 61 z^{96} + 29 z^{97} + 12 z^{98} + 49 z^{99} + O(z^{100}) \\[4pt]
% 71 & 1 + 67 z^{71} + O(z^{100}) \\[4pt]
\hline
\end{array}
\end{equation*}
\caption{ $A_p(z)$ of $\mathcal{L}$ for small primes $p$.}
\label{table: A_p(z) at small primes for pullback of operator 1hat}
\end{table}





% \begin{table}
% \begin{equation*}
% \begin{array}{|c|l|}
% \hline
%  p & A_p(z(q)) \\[2pt]
%  \hline
% 2 & 1 + O(q^{50}) \\[2pt]
% 3 & 1 + O(q^{50}) \\[2pt]
% 5 & 1 + O(q^{50}) \\[2pt]
% 7 & 1 + O(q^{50}) \\[2pt]
% 11 & 1 + 7q^{11} + 9q^{22} + 10q^{33} + 9q^{44} + O(q^{50}) \\[2pt]
% 13 & 1 + 6q^{13} + 10q^{26} + 6q^{39} + O(q^{50}) \\[2pt]
% 17 & 1 + 12q^{17} + 8q^{34} + O(q^{50}) \\[2pt]
% 19 & 1 + 2q^{19} + 16q^{38} + O(q^{50}) \\[2pt]
% 23 & 1 + 22q^{23} + 12q^{24} + 2q^{25} + 3q^{26} + 21q^{27} + 21q^{28} + 7q^{29} + 16q^{30} + 5q^{32} + 4q^{33} + q^{34} + 22q^{35} + 18q^{36}\\[2pt]
%    & + 13q^{37} + 8q^{38} + 12q^{39} + 6q^{40} + 20q^{41} + 6q^{42} + 19q^{43} + 9q^{45} + 16q^{46} + 12q^{48} + 5q^{49} + O(q^{50}) \\[2pt]
% 29 &1 + 13q^{29} + 8q^{30} + 12q^{31} + 18q^{32} + 13q^{33} + 16q^{35} + 18q^{36} + 17q^{37} + 14q^{38} + 12q^{39} + 8q^{40} + 27q^{41} \\[2pt]
%    &+ 18q^{42}+ q^{43} + 6q^{44} + 20q^{45} + 12q^{46} + 14q^{47} + 11q^{48} + 21q^{49} + O(q^{50}) \\[2pt]
% 31 & 1 + 24q^{31} + O(q^{50}) \\[2pt]
% 37 & 1 + 6q^{37} + O(q^{50}) \\[2pt]
% 41 & 1 + 34q^{41} + 40q^{42} + 36q^{43} + 24q^{44} + 30q^{45} + 14q^{46} + 36q^{47} + 24q^{48} + 17q^{49} + O(q^{50}) \\[2pt]
% 43 & 1 + 40q^{43} + 35q^{44} + 20q^{45} + 40q^{46} + 8q^{47} + 24q^{48} + 27q^{49} + O(q^{50}) \\[2pt]
% 47 & 1 + 17q^{47} + 22q^{48} + 31q^{49} + O(q^{50}) \\[2pt]
% \hline
% \end{array}
% \end{equation*}
% \caption{ $A_p(z(q))$ of $\mathcal{L}$ for small primes $p$.}
% \label{table: A_p(z(q)) at small primes for pullback of operator 1hat}
% \end{table}


% For the primes $p\in\{11,13,17,19,31,37\}$, we can see in Table~\ref{table: A_p(z) at small primes for pullback of operator 1hat} that $A_p(z(q))$ contains only $p^{th}$ powers of $q$, so it must be of the form $A_p(z(q))\equiv_pf(q)^p$ for some $f(q)=\mathbb{F}_p\llbracket q\rrbracket$. We check many more primes in an attached Jupyter notebook and conjecture that this happens when $p$ is inert in $\mathbb{Q}\left(\sqrt{-5}\right)$. 

% We compute $A_p(z(q))$ for many primes and use the Chinese remainder theorem to find $f(q)\in(\mathbb{Z}/N\mathbb{Z})\llbracket q\rrbracket$ for large $N\in\mathbb{Z}$. We then use the rational reconstruction algorithm to lift this to $f(q)\in 1+q\mathbb{Z}\llbracket q\rrbracket$ \cite{villegas2007experimental}. As we increase $N$, we find that this lift stabilises to
% \begin{equation}
%     f(q)=1-12500q+128625000q^2-1297013750000q^3+13057738862500000q^4+O(q^5),
% \end{equation}
% which leads to the following conjecture.
% \begin{conj}
% \label{01022026aquamatica}
% Let $p$ be an inert rational prime in $\mathbb{Q}\left(\sqrt{-5}\right)$. The identity
% \begin{equation}
%     A_p(z(q))\equiv_p f(q)^p
% \end{equation}
% holds for some \begin{equation}
%     A_p(z)\equiv_p f_p(z)^p.
% \end{equation}, which is independent of $p$.
% \end{conj}


















% \section{Fourteen examples and AESZ list}
% \def\losu {{G}}

% Let $a_i,\ i=1,2,\ldots,n$ be rational numbers, $0<a_i<1$,   $a=(a_1,a_2,\ldots,a_n)$ and
% $$
% F(a| z):=\   _{n}F_{n-1}(a_1,\ldots,a_n; 1,1,\ldots,1 | z)= \sum_{k=0}^\infty \frac{(a_1)_k\dots(a_n)_k}{k!^n} \, z^k, \quad |z|<1
% $$
%  be the holomorphic solution of the generalized hypergeometric differential equation 
% \begin{equation}
% \label{14fev}
% \theta^n-z(\theta+a_1)(\theta+a_2)\cdots(\theta+a_n)=0
% \end{equation}
% where $(a_i)_k=a_i(a_i+1)(a_i+2)...(a_i+k-1),\,(a_i)_0 = 1,$ is the Pochhammer symbol and $\theta=z\frac{d}{dz}$. 
% The first logarithmic solution in the Frobenius basis around $z=0$ has the form $\losu(a|z)+F(a|z)\log z$,  
% where
% \begin{equation}
% \label{logsol}
% \losu(a| z)=\sum_{k=1}^\infty \frac{(a_1)_k\cdots(a_n)_k}{(k!)^n}\big{[}\sum_{j=1}^ n\sum_{i=0}^{k-1}(\frac{1}{a_j+i}-\frac{1}{1+i})\big{]}z^k.
% \end{equation}
% The mirror map 
% $$
% q(a| z)=:\frac{z}{N}\exp(\dfrac{\losu(a|z)}{F(a| z)} ),
% $$ is a natural generalization of the
% Schwarz function, where $N\in\mathbb{Z}$ is some normalising constant. 
% %This number is uniquely determined by the fact that 

% In this article we are only interested in the case $n=4$ and $a_3=1-a_1$ and $a_4=1-a_2$ and $(a_1,a_2)$ from the Table below:
% \begin{center}
% \begin{tabular}{|c|c|c|}
% \hline
%  $(\alpha_1,\alpha_2)$ & $(d,k)$  & $3\alpha_2-1-\alpha_1$ \\
% \hline
% $(\frac{1}{5},\frac{2}{5})$&$(5,5)$  & + \\
% \hline
% $(\frac{1}{6},\frac{1}{3})$& $(3,4)$& -\\
% \hline
% $(\frac{1}{8},\frac{3}{8})$&$(2,4)$& +\\
% \hline
%  $(\frac{1}{10},\frac{3}{10})$& $(1,3)$& -\\
% \hline
%  $(\frac{1}{4},\frac{1}{3})$& $(6,5)$& -\\
% \hline
% $(\frac{1}{6},\frac{1}{4})$& $(2,3)$& -\\
% \hline
% $(\frac{1}{12},\frac{5}{12})$& $(1,4)$& +\\
% \hline
% $(\frac{1}{4},\frac{1}{2})$&$(8,6)$& +\\
% \hline
% $(\frac{1}{3},\frac{1}{2})$&$(12,7)$& +\\
% \hline
% $(\frac{1}{6},\frac{1}{2})$&$(4,5)$& +\\
% \hline
% $(\frac{1}{3},\frac{1}{3})$&$(9,6)$& -\\
% \hline
% $(\frac{1}{4},\frac{1}{4})$&$(4,4)$& -\\
% \hline
% $(\frac{1}{6},\frac{1}{6})$&$(1,2)$& -\\
% \hline
% $(\frac{1}{2},\frac{1}{2})$&$(16,8)$& + \\ \hline
% \end{tabular}\\
% Data of 14  hypergeometric equations
% \end{center}

% \marginpar{\tiny {\color{red} Mohamed} Definition of a Calabi-Yau equation  and the computer data availeble in github etc..}















\newpage
\bibliography{biblio.bib}
\bibliographystyle{alpha}

\printindex












































































































\end{document}





























\section{Mirror of intersection of four quadrics in $\mathbb{P}^7$}

I work with parameter $\widetilde{z}$ such that the holomorphic period $\varpi_0$ has coefficients in $(2\pi i)^3\mathbb{Z}$ when expanded around the MUM point $\widetilde{z}=0$.

I choose
\begin{align}
    \begin{split}
        t_0 &= \sum_{n=0}^\infty \binom{2n}{n}^4\widetilde{z}^n\\
        t_4 &= z\,t_0^2
    \end{split}
\end{align}

The standard mirror map at $\widetilde{z}=0$ gives
\begin{equation}
    q-64 q^2+1120 q^3-38912 q^4-1536464 q^5+O\left(q^6\right)
\end{equation}
which leads to the expressions
\begin{align}
    \begin{split}
        t_0(q)&=1+16 q+272 q^2+12032 q^3+878864 q^4+78640128 q^5+O\left(q^6\right)\\
   t_4(q) &= q-32 q^2-128 q^3-21504 q^4-1766064 q^5+O\left(q^6\right)\\
    \end{split}
\end{align}

We can find the correct periods with the holomorphic $3$-form
\begin{equation}
    \Omega(\psi)=\frac{1}{2\pi i}\int_{T(P)}\frac{dY_1\wedge dY_2\wedge dY_3\wedge dY_4}{P(Y,\psi)}=\frac{dY_1\wedge dY_2\wedge dY_3}{\partial_{Y_4}P}\Bigg|_{P=0}
\end{equation}
where $T(P)$ is a ``small contour" around $P=0$ and
\begin{equation}
    P(Y,\psi)=\prod_{j=1}^4Y_j-\frac{1}{(2\psi)^4}\prod_{j=1}^4(Y_j^2-1)
\end{equation}
\begin{equation}
    \int_{T^3}\Omega(\psi)=\sum_{n=0}^\infty\binom{2n}{n}^4\left(\frac{1}{2\psi}\right)^{8n}.
\end{equation}

It is more standard to find the mirror of the intersection of four quadrics in $\mathbb{P}^7$ by considering a $(\mathbb{Z}/4\mathbb{Z})^3$ quotient of the family defined by the four polynomials 
\begin{equation}
    P_j=x_j^2+y_j^2-2\psi x_{j+1} y_{j+1}=0
\end{equation}
where $[x_j:y_j]$ are homogeneous coordinates in $\mathbb{P}^7$ and the indices are defined modulo $4$. The holomorphic $3$-form is given by
\begin{equation}
    \Omega(\psi)=\frac{(2\psi)^4}{(2\pi i)^4}\int_{\prod T(P_j)} \frac{dx_0\wedge dy_0\wedge...\wedge dx_3\wedge dy_3}{P_0P_1P_1P_3}
\end{equation}

Moreover according to the computation of the Cartier operator, we have:
\begin{equation}
 A(\psi) = \sum^{\lfloor \frac{p-1}{2} \rfloor}_{i=0} (\frac{(2i)!}{(i!)^2})^4 (-2\psi)^{4(p-1-2i)}.
\end{equation}


\begin{equation*}
\begin{array}{c|l}
p & A_p\\
\hline
 2 & X \\
 3 & X^2+Y \\
 5 & X^4+X^2 Y+Y^2 \\
 7 & X^6+2 X^4 Y+X^2 Y^2+Y^3 \\
 11 & X^{10}+5 X^8 Y+9 X^6 Y^2+5 X^4 Y^3+3 X^2 Y^4+Y^5 \\
 13 & X^{12}+3 X^{10} Y+9 X^8 Y^2+9 X^6 Y^3+X^4 Y^4+X^2 Y^5+Y^6 \\
 17 & X^{16}+16 X^{14} Y+4 X^{12} Y^2+13 X^{10} Y^3+16 X^8 Y^4+13 X^6 Y^5+4 X^4 Y^6+16 X^2
   Y^7+Y^8 \\
 19 & X^{18}+16 X^{16} Y+4 X^{14} Y^2+X^{12} Y^3+4 X^{10} Y^4+17 X^8 Y^5+7 X^6 Y^6+7 X^4 Y^7+7
   X^2 Y^8+Y^9 \\
 23 & X^{22}+16 X^{20} Y+8 X^{18} Y^2+12 X^{16} Y^3+X^{14} Y^4+X^{12} Y^5+3 X^{10} Y^6+4 X^8
   Y^7+18 X^6 Y^8+16 X^4 Y^9+12 X^2 Y^{10}+Y^{11} \\
 29 & X^{28}+16 X^{26} Y+20 X^{24} Y^2+7 X^{22} Y^3+X^{20} Y^4+7 X^{18} Y^5+24 X^{16} Y^6+24
   X^{14} Y^7+20 X^{12} Y^8+25 X^{10} Y^9+23 X^8 Y^{10}+23 X^6 Y^{11}+23 X^4 Y^{12}+25 X^2
   Y^{13}+Y^{14} \\
 31 & X^{30}+16 X^{28} Y+25 X^{26} Y^2+9 X^{24} Y^3+4 X^{22} Y^4+8 X^{20} Y^5+25 X^{18} Y^6+20
   X^{16} Y^7+5 X^{14} Y^8+28 X^{12} Y^9+8 X^{10} Y^{10}+8 X^8 Y^{11}+5 X^6 Y^{12}+14 X^4
   Y^{13}+8 X^2 Y^{14}+Y^{15} \\
 37 & X^{36}+16 X^{34} Y+X^{32} Y^2+12 X^{30} Y^3+34 X^{28} Y^4+33 X^{26} Y^5+X^{24} Y^6+12
   X^{22} Y^7+X^{20} Y^8+16 X^{18} Y^9+9 X^{16} Y^{10}+10 X^{14} Y^{11}+26 X^{12} Y^{12}+26
   X^{10} Y^{13}+9 X^8 Y^{14}+9 X^6 Y^{15}+16 X^4 Y^{16}+10 X^2 Y^{17}+Y^{18} \\
 41 & X^{40}+16 X^{38} Y+25 X^{36} Y^2+18 X^{34} Y^3+31 X^{32} Y^4+25 X^{30} Y^5+23 X^{28} Y^6+31
   X^{26} Y^7+10 X^{24} Y^8+25 X^{22} Y^9+37 X^{20} Y^{10}+40 X^{18} Y^{11}+X^{16} Y^{12}+23
   X^{14} Y^{13}+40 X^{12} Y^{14}+25 X^{10} Y^{15}+25 X^8 Y^{16}+10 X^6 Y^{17}+4 X^4 Y^{18}+10
   X^2 Y^{19}+Y^{20} \\
 43 & X^{42}+16 X^{40} Y+6 X^{38} Y^2+40 X^{36} Y^3+4 X^{34} Y^4+6 X^{32} Y^5+35 X^{30} Y^6+11
   X^{28} Y^7+9 X^{26} Y^8+9 X^{24} Y^9+14 X^{22} Y^{10}+15 X^{20} Y^{11}+14 X^{18} Y^{12}+13
   X^{16} Y^{13}+11 X^{14} Y^{14}+11 X^{12} Y^{15}+10 X^{10} Y^{16}+41 X^8 Y^{17}+6 X^6 Y^{18}+38
   X^4 Y^{19}+4 X^2 Y^{20}+Y^{21} \\
 47 & X^{46}+16 X^{44} Y+27 X^{42} Y^2+12 X^{40} Y^3+3 X^{38} Y^4+2 X^{36} Y^5+18 X^{34}
   Y^6+X^{32} Y^7+7 X^{30} Y^8+8 X^{28} Y^9+X^{26} Y^{10}+28 X^{24} Y^{11}+24 X^{22} Y^{12}+2
   X^{20} Y^{13}+6 X^{18} Y^{14}+24 X^{16} Y^{15}+8 X^{14} Y^{16}+7 X^{12} Y^{17}+14 X^{10}
   Y^{18}+2 X^8 Y^{19}+3 X^6 Y^{20}+4 X^4 Y^{21}+27 X^2 Y^{22}+Y^{23} \\
 53 & X^{52}+16 X^{50} Y+24 X^{48} Y^2+46 X^{46} Y^3+46 X^{44} Y^4+47 X^{42} Y^5+X^{40} Y^6+47
   X^{38} Y^7+42 X^{36} Y^8+47 X^{34} Y^9+16 X^{32} Y^{10}+16 X^{30} Y^{11}+36 X^{28} Y^{12}+44
   X^{26} Y^{13}+X^{24} Y^{14}+36 X^{22} Y^{15}+X^{20} Y^{16}+16 X^{18} Y^{17}+44 X^{16}
   Y^{18}+36 X^{14} Y^{19}+44 X^{12} Y^{20}+28 X^{10} Y^{21}+49 X^8 Y^{22}+47 X^6 Y^{23}+46 X^4
   Y^{24}+46 X^2 Y^{25}+Y^{26} \\
 59 & X^{58}+16 X^{56} Y+57 X^{54} Y^2+51 X^{52} Y^3+9 X^{50} Y^4+46 X^{48} Y^5+51 X^{46} Y^6+29
   X^{44} Y^7+25 X^{42} Y^8+20 X^{40} Y^9+28 X^{38} Y^{10}+53 X^{36} Y^{11}+12 X^{34} Y^{12}+17
   X^{32} Y^{13}+26 X^{30} Y^{14}+48 X^{28} Y^{15}+5 X^{26} Y^{16}+27 X^{24} Y^{17}+28 X^{22}
   Y^{18}+27 X^{20} Y^{19}+19 X^{18} Y^{20}+X^{16} Y^{21}+51 X^{14} Y^{22}+12 X^{12} Y^{23}+12
   X^{10} Y^{24}+49 X^8 Y^{25}+9 X^6 Y^{26}+15 X^4 Y^{27}+26 X^2 Y^{28}+Y^{29} \\
 61 & X^{60}+16 X^{58} Y+15 X^{56} Y^2+58 X^{54} Y^3+34 X^{52} Y^4+9 X^{50} Y^5+34 X^{48} Y^6+22
   X^{46} Y^7+X^{44} Y^8+20 X^{42} Y^9+13 X^{40} Y^{10}+20 X^{38} Y^{11}+25 X^{36} Y^{12}+16
   X^{34} Y^{13}+X^{32} Y^{14}+57 X^{30} Y^{15}+22 X^{28} Y^{16}+58 X^{26} Y^{17}+57 X^{24}
   Y^{18}+15 X^{22} Y^{19}+X^{20} Y^{20}+X^{18} Y^{21}+56 X^{16} Y^{22}+20 X^{14} Y^{23}+9 X^{12}
   Y^{24}+56 X^{10} Y^{25}+25 X^8 Y^{26}+34 X^6 Y^{27}+42 X^4 Y^{28}+34 X^2 Y^{29}+Y^{30} \\
 67 & X^{66}+16 X^{64} Y+23 X^{62} Y^2+4 X^{60} Y^3+14 X^{58} Y^4+10 X^{56} Y^5+25 X^{54} Y^6+40
   X^{52} Y^7+23 X^{50} Y^8+24 X^{48} Y^9+37 X^{46} Y^{10}+35 X^{44} Y^{11}+60 X^{42} Y^{12}+35
   X^{40} Y^{13}+40 X^{38} Y^{14}+62 X^{36} Y^{15}+56 X^{34} Y^{16}+65 X^{32} Y^{17}+64 X^{30}
   Y^{18}+39 X^{28} Y^{19}+23 X^{26} Y^{20}+21 X^{24} Y^{21}+22 X^{22} Y^{22}+22 X^{20} Y^{23}+25
   X^{18} Y^{24}+33 X^{16} Y^{25}+35 X^{14} Y^{26}+X^{12} Y^{27}+4 X^{10} Y^{28}+56 X^8 Y^{29}+26
   X^6 Y^{30}+21 X^4 Y^{31}+15 X^2 Y^{32}+Y^{33} \\
 71 & X^{70}+16 X^{68} Y+18 X^{66} Y^2+37 X^{64} Y^3+X^{62} Y^4+48 X^{60} Y^5+X^{58} Y^6+64
   X^{56} Y^7+36 X^{54} Y^8+2 X^{52} Y^9+5 X^{50} Y^{10}+64 X^{48} Y^{11}+12 X^{46} Y^{12}+30
   X^{44} Y^{13}+8 X^{42} Y^{14}+5 X^{40} Y^{15}+X^{38} Y^{16}+25 X^{36} Y^{17}+10 X^{34}
   Y^{18}+58 X^{32} Y^{19}+18 X^{30} Y^{20}+58 X^{28} Y^{21}+49 X^{26} Y^{22}+2 X^{24} Y^{23}+32
   X^{22} Y^{24}+43 X^{20} Y^{25}+9 X^{18} Y^{26}+60 X^{16} Y^{27}+36 X^{14} Y^{28}+15 X^{12}
   Y^{29}+30 X^{10} Y^{30}+64 X^8 Y^{31}+4 X^6 Y^{32}+2 X^4 Y^{33}+29 X^2 Y^{34}+Y^{35} 
\end{array}
\end{equation*}




\section{Hypersurfaces in weighted projective space $\mathbb{P}^4(k_0,k_1,k_2,k_3,k_4)$}

These examples can be found in \cite{Morrison:1991cd}. They are given by the hypersurfaces in the weighted projective space $\mathbb{P}^4(k_0,k_1,k_2,k_3,k_4)$.

% $$
% \begin{array}{c|c|c}
%    AESZ \# & \mathbb{P}[k_i] &  P(x,\psi)\\
%    \hline
% 1 & (1,1,1,1,1) & M_5x_0^5 + x_1^5 + x_2^5+x_3^5 + x_4^5 -M_1x_0x_1x_2x_3x_4\\
% 2 & (1,1,1,2,5) & M_{10}x_0^2 + x_1^5 + x_2^{10} + x_3^{10} + x_4^{10} - M_1 x_0x_1x_2x_3x_4\\
% 7 & (1,1,1,1,4) & M_8x_0^2 + x_1^8 + x_2^8 + x_3^8 + x_4^8 - M_1x_0x_1x_2x_3x_4\\
% 8 &   (1,1,1,1,2) & M_6x_0^6 + x_1^6 + x_2^6 + x_3^3 + x_4^6 - M_1x_0x_1x_2x_3x_4\\
% \end{array}
% $$
$$
\begin{array}{c|c|c}
   AESZ \# & \mathbb{P}[k_i] &  P(x,\psi)\\
   \hline
1 & (1,1,1,1,1) & M_5x_0^5 + x_1^5 + x_2^5+x_3^5 + x_4^5 -M_1x_0x_1x_2x_3x_4\\
2 & (5,2,1,1,1) & x_0^2 + x_1^5 + M_{10}x_2^{10} + x_3^{10} + x_4^{10} - M_1 x_0x_1x_2x_3x_4\\
7 & (4,1,1,1,1) & x_0^2 + M_8x_1^8 + x_2^8 + x_3^8 + x_4^8 - M_1x_0x_1x_2x_3x_4\\
8 &   (1,1,1,2,1) & M_6x_0^6 + x_1^6 + x_2^6 + x_3^3 + x_4^6 - M_1x_0x_1x_2x_3x_4\\
\end{array}
$$

The differential $3$-form is obtained by 
$$
\alpha=\frac{dx_1\wedge dx_2\wedge dx_3\wedge dx_4}{df}
$$
where $f$ is the equation of mirror $CY3$  with $x_0=1$. For the sake of computation we can take $M_1=x$ (Mohammed change X to x and $Y$ to $y$)
and $M_{i}=y,\ i=5,10,6,8$.

According to Katz's formula, we could derive the following formulas.
\begin{itemize}
    \item{(AESZ 2)}
    \begin{equation}
A(x, y) = \sum_{i=0}^{\lfloor \frac{p-1}{10} \rfloor} \frac{(p-1)!}{((5i)!)(i!)^3((2i)!)(p-1-10i)!} y^{i} x^{p-1-10i}. 
\end{equation}
    \item{(AESZ 7)}
    \begin{equation}
A(x, y) = \sum_{i=0}^{\lfloor \frac{p-1}{8} \rfloor} \frac{(p-1)!}{(i!)^4((4i)!)(p-1-8i)!} y^{i} x^{p-1-8i}. 
\end{equation}
    \item{(AESZ 8)}
    \begin{equation}
A(x, y) = \sum_{i=0}^{\lfloor \frac{p-1}{6} \rfloor} \frac{(p-1)!}{(i!)^4((2i)!)(p-1-6i)!} y^i x^{p-1-6i}. 
\end{equation}
\end{itemize}

For (AESZ 2),
\[
\begin{aligned}
p=2: &\quad x \\
p=3: &\quad x^2 \\
p=5: &\quad x^4 \\
p=7: &\quad x^6 \\
p=11: &\quad x^{10} + 6y \\
p=13: &\quad x^{12} + x^2y\\
p=17: &\quad x^{16} + 7x^{6}y \\
p=19: &\quad x^{18} + 15x^{8}y  \\
p=23: &\quad x^{22}+9x^{12}y+18x^2y^2\\
p=29: &\quad x^{28}+11x^{18}y+24x^{8}y^2 \\
\end{aligned}
\]

For (AESZ 7),
\[
\begin{aligned}
p=2: &\quad x \\
p=3: &\quad x^2 \\
p=5: &\quad x^4 \\
p=7: &\quad x^6 \\
p=11: &\quad x^{10}+ 8x^2y^4\\
p=13: &\quad x^{12}+ 3x^4y^4\\
p=17: &\quad x^{16}+14x^8y^4+4y^8 \\
p=19: &\quad x^{18} + 8x^{10}y^4+8x^2y^8\\
p=23: &\quad x^{22}+x^{14}y^4+8x^6y^8 \\
p=29: &\quad x^{28}+27x^{20}y^4+18x^{12}y^8+6x^4y^{12}\\
\end{aligned}
\]

For (AESZ 8),
\[
\begin{aligned}
p=2: &\quad x \\
p=3: &\quad x^2 \\
p=5: &\quad x^4 \\
p=7: &\quad x^6 + 3y \\
p=11: &\quad x^{10} + 8x^4y \\
p=13: &\quad x^{12} + 9x^6y + 11y^2 \\
p=17: &\quad x^{16} + 3x^{10}y + 8x^4y^2 \\
p=19: &\quad x^{18} + 18x^{12}y + 12x^6y^2 + 12y^3 \\
p=23: &\quad x^{22}+15x^{16}y+18x^{10}y^2+16x^4y^3 \\
p=29: &\quad x^{28}+12x^{22}y+23x^{16}y^2+9x^{10}y^3+12x^4y^4 \\
\end{aligned}
\]


% \subsection{AESZ $8$}


% $$
% \begin{array}{c|l}
% p & A_p\\
% \hline
%  2 & x \\
%  3 & x^2 \\
%  5 & x^4 \\
%  7 & x^6+3 y \\
%  11 & x^{10}+8 x^4 y \\
%  13 & x^{12}+9 x^6 y+11 y^2 \\
%  17 & x^{16}+3 x^{10} y+8 x^4 y^2 \\
%  19 & x^{18}+18 x^{12} y+12 x^6 y^2+12 y^3 \\
%  23 & x^{22}+15 x^{16} y+18 x^{10} y^2+16 x^4 y^3 \\
%  29 & x^{28}+12 x^{22} y+23 x^{16} y^2+9 x^{10} y^3+12 x^4 y^4 \\
%  31 & x^{30}+19 x^{24} y+22 x^{18} y^2+3 x^{12} y^3+22 x^6 y^4+29 y^5 \\
%  37 & x^{36}+27 x^{30} y+19 x^{24} y^2+4 x^{18} y^3+34 x^{12} y^4+31 x^6 y^5+11 y^6 \\
%  41 & x^{40}+32 x^{34} y+16 x^{28} y^2+12 x^{22} y^3+15 x^{16} y^4+33 x^{10} y^5+36 x^4 y^6 \\
%  43 & x^{42}+16 x^{36} y+13 x^{30} y^2+35 x^{24} y^3+26 x^{18} y^4+13 x^{12} y^5+41 x^6 y^6+39 y^7 \\
%  47 & x^{46}+31 x^{40} y+20 x^{34} y^2+8 x^{28} y^3+17 x^{22} y^4+25 x^{16} y^5+14 x^{10} y^6+19 x^4 y^7 \\
%  53 & x^{52}+42 x^{46} y+45 x^{40} y^2+10 x^{34} y^3+40 x^{28} y^4+26 x^{22} y^5+41 x^{16} y^6+37 x^{10} y^7+17 x^4 y^8 \\
%  59 & x^{58}+6 x^{52} y+22 x^{46} y^2+27 x^{40} y^3+38 x^{34} y^4+52 x^{28} y^5+16 x^{22} y^6+43 x^{16} y^7+26 x^{10} y^8+19 x^4 y^9 \\
%  61 & x^{60}+55 x^{54} y+11 x^{48} y^2+20 x^{42} y^3+59 x^{36} y^4+3 x^{30} y^5+51 x^{24} y^6+36 x^{18} y^7+57 x^{12} y^8+45 x^6 y^9+14 y^{10}
%    \\
%  67 & x^{66}+25 x^{60} y+61 x^{54} y^2+6 x^{48} y^3+45 x^{42} y^4+33 x^{36} y^5+3 x^{30} y^6+42 x^{24} y^7+36 x^{18} y^8+2 x^{12} y^9+17 x^6
%    y^{10}+2 y^{11} \\
%  71 & x^{70}+5 x^{64} y+x^{58} y^2+65 x^{52} y^3+13 x^{46} y^4+46 x^{40} y^5+x^{34} y^6+46 x^{28} y^7+9 x^{22} y^8+10 x^{16} y^9+70 x^{10}
%    y^{10}+4 x^4 y^{11} \\
% \end{array}
% $$


\subsection{Guess for general Cartier matrix}

The four examples of Calabi-Yau hypersurfaces in weighted projective space all have the same form. The enhanced moduli space is parametrised by $x$ and $y$. We assume that
\begin{align}
\label{eq: q-expansions for x and y}
    x=&\varpi_0\\
    y=&z\varpi_0^\beta
\end{align}
for some $\beta$, where $z$ is a coordinate on the non-enhanced moduli space such that the holomorphic period
\begin{equation}
    \varpi_0(z)=\int_{T^3}\Omega(z)\in(2\pi i)^3\mathbb{Z} \llbracket z\rrbracket
\end{equation}
for some three torus $T^3$. There is a $\mathbb{C}^\times$ action, such that $\Omega\mapsto k\,\Omega$ for some $k\in\mathbb{C}^\times$. The coordinates $x$ and $y$ must also transform with some weight, which we find is $(x,y)\mapsto(kx,\,k^\beta y)$. This fixes the powers of $\varpi_0$, which appear on the right hand side of \eqref{eq: q-expansions for x and y}. 

\ul{ It is still unclear to me why we choose $\varpi_0$. The mirror map $z(q)$ has a concrete geometric meaning. I do not know what it has to do with the enhanced moduli space.}


We find that 
\begin{equation}
    A_p(x,y)=x^{p-1}\sum_{i=0}^{\lfloor\frac{p-1}{\beta}\rfloor}a_i\binom{p-1}{\beta i}(-x)^{-\beta i}y^i
\end{equation}
where $\beta\in\mathbb{Z}$ and $a_i\in\mathbb{Z}$ are the coefficients that appear in the holomorphic period.

We can check with a computer that, for all fourteen hypergeometric Calabi-Yau operators and a number of non-hypergeometric operators, there is a $\beta\in\mathbb{Z}$ such that
\begin{equation}
    A_p(x(q),y(q))=1+\mathcal{O}(q^{100}).
\end{equation}
for the first 20 primes.

\begin{rem}
In the hypergeometric cases, the integer $\beta$ seems to be related to the order of the orbifold point at $z=\infty$. We observe that in many hypergeometric and non-hypergeometric examples, there is a maximal $\beta_{max}$ such that the $q$ series $A_p(x(q),y(q))=1$. Moreover, any $\beta<\beta_{max}$ will lead to $A_p(x(q),y(q))=1$, when $x(q)$ and $y(q)$ are defined as in \eqref{eq: q-expansions for x and y}.
\end{rem}

\begin{rem}
It has been observed that there is sometimes an integer $N_{max}$ such that
\begin{equation}
    \left(\frac{1}{z}q(z)\right)^{\frac{1}{N_{max}}}\in\mathbb{Z}\llbracket z\rrbracket
\end{equation}

Perhaps $N_{max}=\beta_{\max}$.
\end{rem}

AESZ34 points out two problems with this. It seems that $N\in\{2,4,8\}$ give integer power series this operator. However, $\beta_{max}=1$ for $A_p$.

\textbf{UPDATE:} It seems that varying $\beta$ in some range $1\leq \beta\leq\beta_{max}$ does not change the polynomial $A_p$. I've checked this for the 14 hypergeometric CY operators for the first 30 primes.

The reason for this behaviour seems to be that many of the integer coefficients appearing in the fundamental period 
\begin{equation}
    \int_{\delta_{0,z}}\alpha_z
\end{equation}
vanish mod $p$. A related observation can be found in section 2.6 of \cite{Candelas:2000fq}.

\textbf{Warning:}
One must be careful with the truncated period in Conjecture XXXXX. We have conjectured that
\begin{equation}
    A(z(q))=\left(\sum_{n=0}^\infty c_n z(q)^n\right)^{p-1} \left(\sum_{n=0}^{p-1}c_nz(q)^n\right) \equiv_p1
\end{equation}
as a formal $q$-expansion with coefficients in $\mathbb{F}_p$, where $c_n\in\mathbb{Z}$ are the coefficients of the expansion of the fundamental period around the MUM point $z=0$. From this, one might be tempted to conclude that
\begin{equation}
    \left(\sum_{n=0}^\infty c_n z(q)^n\right)^{p}\equiv_p 1.
\end{equation}

However, this is not true. As an example, consider the operator AESZ 1. The mirror map is given by
\begin{equation}
    z(q)=q-770 q^2+171525 q^3-81623000 q^4-35423171250 q^5 + \mathcal{O}(q^{6})
\end{equation}
and we can confirm for the first 30 primes that
\begin{equation}
    \left(\sum_{n=0}^\infty \frac{(5n)!}{(n!)^5} z(q)^n\right)^{p-1} \left(\sum_{n=0}^{p-1}a_nz(q)^n\right) \equiv_p 1 + \mathcal{O}(q^{201})
\end{equation}
However, we find that
\begin{equation}
    \left(\sum_{n=0}^{7-1}\frac{(5n)!}{(n!)^5} z(q)^n\right)^{7}\equiv_p1+q^7+4 q^{21}+3 q^{28}+2 q^{35}+5 q^{42}+4 q^{49}+4 q^{56}+O\left(q^{63}\right)
\end{equation}




\section{A non-hypergeometric example}

For a non-hypergeometric family, we consider $P=0$, where
\begin{equation}
   P= X_1X_2X_3X_4X_5\left[1-\varphi\left(X_1+X_2+X_3+X_4+X_5\right)\left(\frac{1}{X_1}+\frac{1}{X_2}+\frac{1}{X_3}+\frac{1}{X_4}+\frac{1}{X_5}\right)\right]
\end{equation}
We should strictly speaking work on a $\mathbb{Z}/5\mathbb{Z}$ or $\mathbb{Z}/10\mathbb{Z}$ quotient of the variety defined by $P=0$.

On the $X_1=1$ coordinate patch, the holomorphic threeform is given by
\begin{equation}
    \Omega(\varphi)=\frac{1}{2\pi i}\int_{T(P)} \frac{dX_2\wedge dX_3\wedge dX_4\wedge dX_5}{P|_{X_1=1}}
\end{equation}

We can compute the integral 
\begin{equation}
    \int_{T^3}\Omega(\varphi)=\frac{1}{2\pi i}\int_{\prod_{i=2}^5T(X_i)} \frac{dX_2\wedge dX_3\wedge dX_4\wedge dX_5}{P|_{X_1=1}}=(2\pi i)^3\sum_{n=0}^\infty a_n \varphi^n
\end{equation}
where
\begin{equation}a_n=\sum_{p+q+r+s+t=n}\left(\frac{n!}{p!q!r!s!t!}\right)^2 .
\end{equation}
This period is annihilated by the Picard-Fuchs operator AESZ 34, which has the Riemann symbol
\begin{equation}
    \mathcal{P}~\left\{~\begin{array}{c c c c c}  0 & \frac{1}{25} & \frac{1}{9} & 1 & \infty\\[5pt]
    \hline
    \\[-5pt]
     0 & 0 & 0 & 0 & 1\\
     0 & 1 & 1 & 1 & 1\\
     0 & 1 & 1 & 1 & 2\\
     0 & 2 & 2 & 2 & 2
    \end{array}~~; ~\varphi\right\}~.
\end{equation}
In other words, there is a MUM points at $\varphi=0$, conifold points at $\varphi\in \left\{\frac{1}{25},\frac{1}{9},1\right\}$, and a $K$-point at $\varphi=\infty$. 

On the $X_1=1$ coordinate patch, we can define a two parameter family of threefolds $X_{(\psi,\varphi)}$, with the defining equation
\begin{equation}
\label{eq: two parameter HV threefold}
    \psi X_2X_3X_4X_5\left[1-\varphi\left(\psi+X_2+X_3+X_4+X_5\right)\left(\frac{1}{\psi}+\frac{1}{X_2}+\frac{1}{X_3}+\frac{1}{X_4}+\frac{1}{X_5}\right)\right].
\end{equation}

Note that there is a $\mathbb{C}^\times$ action on $(\psi,\varphi)$. Given $k\in\mathbb{C}^\times$
\begin{equation}
    (\psi,\varphi)\mapsto \left(k\,\psi,\,\varphi\right)
\end{equation}
By combining this $\mathbb{C}^\times$ action with the change of coordinates
\begin{equation}
    X_j\mapsto k X_j
\end{equation}
for $j\in\{2,3,4,5\}$, we see that equation \eqref{eq: two parameter HV threefold} simply picks up an overall factor of $k^5$ and
\begin{equation}
    \Omega\mapsto\frac{1}{k}\Omega
\end{equation}

It follows that $(\psi,\,\varphi)$ parametrise inequivalent pairs $\left( X_{(\psi,\varphi)},\Omega\right)$.



\begin{rem}
  The formula for the Cartier operator seems to be:
\[
\psi^{p-1}(1-5\varphi)^{p-1}.
\]  
\end{rem}






\section{Calabi-Yau Fourfolds}

The Picard-Fuchs equation for the simplest Calabi-Yau fourfolds can be found in \cite{aboBizet:2014ovf}.











\section{Cartier-Manin matrices and Hasse-Witt matrices for hypersurfaces}
\subsection{General construction}
\textbf{The Hass-Witt matrix}
\

We recall the definition of the Hasse-Witt operator for a family of scheme as follows:
\

Give a smooth family of schemes $f: X \to S$, where $S$ is a smooth $T$-scheme for any base scheme $T$ with characteristic $p > 0$. We let $D$ be a union of divisors $D_i$ in $X$, each of them is smooth over $S$ and which have normal crossings relative to $S$.
Then it is well-known that we have two spectral sequences:
\begin{itemize}
    \item{(The Hodge to de Rham spectral sequence)}
    \[
    E_1^{p,q} = R^qf_*(\Omega^p_{X/S}(\log D)) \Rightarrow \mathbb{R}^{p+q}f_*(\Omega^{\bullet}_{X/S}(\log D)).
    \]
    \item{(The conjugate spectral sequence)}
    \[
    _{con} E_2^{p,q} = R^qf_*(\mathcal{H}^q(\Omega^{\bullet}_{X/S}(\log D))) \Rightarrow \mathbb{R}^{p+q}f_*(\Omega^{\bullet}_{X/S}(\log D)).
    \]
\end{itemize}
Via the Cartier isomorphism, we have an isomorphism:
\[
_{con} E_2^{a,b} = R^af_*(\mathcal{H}^b(\Omega^{\bullet}_{X/S}(\log D))) \xrightarrow{C}  R^af_*(\sigma^*(\Omega^{b}_{X/S}(\log D))),
\]
where $\sigma: X^{(p)} \to X$ is the base change of the Frobenius morphism on $S$.
Based on this fact together with some suitable assumption (See \cite[Proposition 2.3.2]{Katz1972}), the Hodge to de Rham spectral sequence degenerate at $E_1$ if and only if the conjugate spectral sequence degenerate at $E_2$. Under these assumptions, Katz defined the Hasse-Witt operator as
\begin{equation}
\mathrm{HW}: F^*_{\mathrm{abs}}R^nf_*(\mathcal{O}_X) \cong _{con} E_2^{n,0} \hookrightarrow \mathbb{R}f_*(\Omega^{\bullet}_{X/S}(\log D)) \to E_1^{0,n} = R^nf_*(\mathcal{O}_X),   
\end{equation}
where the first isomorphism is induced by the inverse of the Cartier isomorphism.

In the case of the family of hypersurfaces $X$ in $\mathbb{P}^{n+1}_S$, Katz \cite[Section 2.3.7]{Katz1972} provided an algorithm for the computation of the Hasse-Witt matrix. The rough idea is as follows:
\begin{itemize}
    \item 
Assume that $X$ is defined by a homogeneous polynomial $H$ of degree $d$ with coefficients in $\Gamma(S, \mathcal{O}_S)$. The short exact sequence
\[
0 \to \mathcal{O}_{\mathbb{P}}(-d) \to \mathcal{O}_{\mathbb{P}} \to \mathcal{O}_X \to 0
\]
induces an isomorphism
\[
R^n f_*(\mathcal{O}_X) \cong R^{n+1} \pi_*(\mathcal{O}_{\mathbb{P}}(-d)),
\]
where $\pi: \mathbb{P}_S^{n+1}\to S$. Since the right hand of the above isomorphism is computable, which admits 
\begin{equation}
X^W = X_0^{W_0} \cdots X^{W_{n+1}}_{n+1}, \sum W_i = -d, W_i < 0, \forall i
\end{equation}
as base, then we may view them as a basis of $R^n f_*(\mathcal{O}_X)$, which is denoted by $m(W)$ in the notation of Katz.
\item 
With the above chosen basis, which we rewrite them as
\[
\{ m(-W) \mid  \sum W_i = d, W_i > 0, \forall i\},
\]
Katz showed that the Hasse-Witt matrix is given by
\begin{equation}
\mathrm{HW}(F^*_{\mathrm{abs}}(m(-W))) = \sum A_{pW-V} (m(-V))
\end{equation}
where $A_U$ is the coefficient of the following polynomial
\[
H^{p-1} = \sum A_U X^U.
\]
\end{itemize}
\textbf{The Cartier-Manin matrix for the Cartier operator}
\

Under the Serre duality between $R^nf_*(\mathcal{O}_X)$ and $f_*(\Omega_{X/S}^n)$, the bases $m(-W)$ ($\sum W_i = d, W_i > 0 \ \forall i$) is corresponding to the bases
\begin{equation}
\omega(W) = \frac{x_1^{W_1-1} \cdots x_{n+1}^{W_{n+1}-1}dx_1 \wedge \cdots \wedge dx_n}{\frac{\partial h}{\partial x_{n+1}}},
\end{equation}
where $x_i = X_i/X_0$ and $h$ is the affine version of $H$. See \cite[Section 2.3.7]{Katz1972} for the details. With the chosen basis $\omega(W)$, the Cartier operator
\begin{equation}
C: f_*(\Omega_{X/S}^n) \to f_*(\Omega_{X/S}^n)
\end{equation}
determines a matrix, which is called the Cartier-Manin matrix. The concrete formula has been shown in \cite[Corollary 1]{Miller72}. In fact, the Cartier-Manin matrix $\tilde{A}_{U,V}$ of the Cartier operator with respect to the above basis $\omega(W)$ is given by the coefficient of $X^V$ in $\psi(H^{p-1}X^U)$, where $\psi$ is the $k$-linear operator on the polynomial ring $k[x_0^{u_1}\cdots x_{n+1}^{u_{n+1}} = x^U]$, which satisfies that
\[
\psi(x^U) = 
\begin{cases}
x^V, & \mathrm{if} \ U=pV; \\
0, & \mathrm{otherwise}.
\end{cases}
\]
\begin{prop}
The Hasse-Witt matrix $A_{U,V}$ is $(\tilde{A}_{U,V})^p$.     
\end{prop}
(This needs to be proved!)
\begin{rem}
This proposition is a generalization of the relation between the Hasse-Witt matrix and the Cartier-Manin matrix in the curve case to the hypersurfaces.
\end{rem}
\begin{exam} \label{basicexample}
Assume the degree $d$ of the hypersurface is $n+2$ and $W = \underbrace{(1, \cdots, 1)}_d$. Then the index $W$ is corresponding to a closed $n$-form $\omega(W) = \frac{dx_1 \wedge \cdots \wedge dx_n}{\frac{\partial h}{\partial x_{n+1}}}$
in the affine coordinates. Then 
\begin{equation}
C(\omega(W)) = (x_1\cdots x_{n+1})^{-1}\psi(h^{p-1}x_1\cdots x_{n+1}) \omega(W).
\end{equation}
\end{exam}

\begin{rem}
Both Katz and Miller's results could be generalized to the case of hypersurfaces in a weighted projective spaces directly. 
\end{rem}





\section{Some thoughts}

\begin{rem}   %%%%%%%%%%%%% Jin 11.16
My feeling is as follows: Based on Proposition 2.3.6.3 in \cite{Katz1972}, we may prove that the Picard-Fuchs equation satisfied by the Hasse-Witt invariant is the same as the Picard-Fuchs equation satisfied by some period. 

Let me give another example. Consider the Weierstrass family
$y^2 = 4x^3 - t_2x -t_3$ over a scheme $S$ of characteristic $p$ together with the holomorphic form $\frac{dx}{y}$. Katz's proposition 2.3.6.3 says that: If the vector field $v$ on $S$ satisfies 
\[
\nabla_{v} (\frac{xdx}{y}) = 0,
\]
then $v(A(t_2, t_3)) = 0$, where $\nabla$ is the Gauss-Manin connection on the family. On the other hand, we know that 
\[
\nabla_{\sf R} (\frac{xdx}{y}) = 0 
\]\marginpar{This is valid only for $\frac{(x+t_1)dx}{y}$. Is there a derivation whose composition of Gauss-Manin connection annihilates $\frac{xdx}{y}-\frac{B(t_2,t_3)dx}{12y}$?}
where $\sf R$ is the Ramanujan vector field. Hence $\mathsf{R}(A(t_2, t_3)) \equiv_p 0$. Using the reduction map $\mathbb{F}_p[t_i] \to \mathbb{F}_p[[q]]$ together with the commutative diagram:
\begin{center}
\begin{tikzcd}
& \mathbb{F}_p[t_i] \arrow[r, "\sf R"] \arrow[d]
& \mathbb{F}_p[t_i] \arrow[d] \\
& \mathbb{F}_p[[q]] \arrow[r, "-q\frac{\partial}{\partial q}"]
& \mathbb{F}_p[[q]]
\end{tikzcd}
\end{center}
we get that:
\[
-q\frac{\partial}{\partial q}(A(t_2(q), t_3(q))) = 0,
\]
and hence $A(t_2(q), t_3(q))$ is a constant mod $p$. 

\end{rem}

\section{Elliptic curves}
In the case of elliptic curves, it turns out that $\Ts=\spec(\Z[\frac{1}{\Delta},t_2,t_3]),\ \ \Delta:=6(27t_3^2-t_2^3)$ and the universal family of elliptic curves over $\Ts$ is given by the classical Weierstrass format $y^2=4x^3-t_2x-t_3$ with $\alpha=\frac{dx}{y}$.  The above statement in the case of elliptic curves is a classical result due to Swinnerton-Dyer, Deligne, Serre. 

\begin{rem}
We let $E$ be a family of elliptic curve over $R$, where $R$ is an algebra over $\mathbb{F}_p$. The absolute Frobenius map $F_{\mathrm{abs}}$ is an additive p-linear endomorphism of $\mathcal{O}_E$ and it induces a p-linear endomorphism of $H^1(E, \mathcal{O}_E)$. We let $\alpha$ be the base of $H^0(E, \Omega^1_{E/R})$ whose dual element $\eta$ is a base of $H^1(E, \mathcal{O}_E)$ under the Serre duality. Then Katz defined $A(E, \alpha) \in R$ as
\[
F_{\mathrm{abs}}^*(\eta) = A(E, \alpha) \eta.
\]
There is another equivalent definition of $A(E, \alpha) \in R$ in \cite[Section 12.4, (12.4.1.4)]{KatzMazur1985}, which is 
\[
C(\alpha) = A(E, \alpha)^{\frac{1}{p}} \alpha,
\]
where $C$ is the Cartier operator. There seems no proof about the equivalence between these two definition in \cite{KatzMazur1985}. We remark that one way to show these two definitions are equivalent is based on the concrete computation. In fact, when $(E,\alpha)$, defined over a perfect field, is given by the Weierstrass equation $y^2=f(x)$ and $\alpha = \frac{dx}{y}$, both of these two definition gives us that $A(E, \alpha)$ is the coefficient of $x^{p-1}$ in $f(x)^{\frac{p-1}{2}}$.  
\

There is a generalization of $A(E, \alpha)$ for any smooth curve $X$ over a perfect field $k$. In \cite{AchterHowe}, the authors suggest that we should make a distinction between  $A(E, \alpha)$ and $ A(E, \alpha)^{\frac{1}{p}}$:
`` If you are working with the Cartier operator on differentials, refer to the matrix representation as the Cartier–Manin matrix; if you are working with the Frobenius operator on $H^1(X, \mathcal{O}_X)$, refer to the matrix representation as the Hasse–Witt matrix." In particular, when $k$ is algebraically closed, Serre showed that the Cartier operator and the Frobenius operator is adjoint under the Serre duality in the following sense:
\[
(\omega, F_{\mathrm{abs}}^*(\eta)) = (C(\omega), \eta)^{\sigma}
\]
where the bracket is the Serre duality and $\sigma$ is the Frobenius map on $k$. See the comments in \cite[Section 2.4]{AchterHowe}. As a corollary, one may derive the relation between the Cartier–Manin matrix and the Hasse–Witt matrix for smooth projective curves over an algebraically closed field.
\end{rem}

\begin{rem}
Let us give a summary about Conjecture \ref{02112025huairou} for elliptic curves. The main reference is Katz's paper \cite[Section 2.0, 2.1]{Katz1972}. 
\begin{itemize}
\item 
Using the above definition of $A(E, \omega)$, one may check directly that $A(E, \omega)$ is a modular form of level one and weight $p-1$ defined over $\mathbb{F}_p$ in the sense of Katz \cite[Section 1.1]{Katz1972}.
\item 
Next we want to see $A$ is holomorphic at $\infty$. By definition we need to show the q-expansion of $A$, which is
\[
A(\mathrm{Tate}(q), \omega_{can}) \in \mathbb{F}_p((q)),
\]
should lie in the subring $\mathbb{F}_p[[q]]$. The key is that the Tate curve over $\mathbb{F}_p((q))$ is the restriction of a plane curve $C$ over $\mathbb{F}_p[[q]]$ and so is $\omega_{can}$ and $\eta_{can}$. Then $A(\mathrm{Tate}(q), \omega_{can})$ is just the matrix $F_{\mathrm{abs}}^*$ on $H^1(C, \mathcal{O}_C)$.
\item 
Katz computed $A(\mathrm{Tate}(q), \omega_{can}) = 1$ by the following method. Note that $H^1(E, \mathcal{O}_E)$ can be identified with the $R$-module of all translation-invariant derivations of $E/R$. Then the Frobenius action on $H^1(E, \mathcal{O}_E)$ is exactly taking the p-th iteration of an invariant derivation. On the completion of the Tate curve along its identity section, one could choose a local coordinate $t$ such that $\omega_{can} = \frac{dt}{1+t}$. Then its dual is the invariant derivation $D(t) = 1+t$. Hence
\[
D^p(1+t) = \cdots = D(1+t)=1+t. 
\]
and then $F_{\mathrm{abs}}^*(\eta_{can}) = \eta_{can}$, which implied that $A(\mathrm{Tate}(q), \omega_{can})=1$. As a corollary, we get 
\[
A \equiv E_{p-1} \mod p,
\]
since these two modular forms over $\mathbb{F}_p$ have the same weight and the same q-expansion when $p\geq 5$.

Alternatively, under the canonical isomorphism \cite[Section 8.8]{KatzMazur1985}(Here we need details.)
\[
\phi_{can}: \widehat{\mathrm{Tate}(q)} \cong \widehat{\mathbb{G}_m},
\]
$\omega_{can} = \phi_{can}^*(\frac{dx}{x})$, where $x$ is the standard multiplication coordinate on $\widehat{\mathbb{G}_m}$. Then the pull back of the derivation $x\frac{\partial}{\partial x}$ of $\widehat{\mathbb{G}_m}$ is corresponding to the dual element of $\omega_{can}$. Since
\[
(x\frac{\partial}{\partial x})^p = x\frac{\partial}{\partial x},
\]
then we get $A(\mathrm{Tate}(q), \omega_{can})=1$. This is the proof given in \cite[Theorem 12.4.2]{KatzMazur1985}.
\end{itemize}
\end{rem}
According to Katz's proof, I'm wondering that:
\begin{itemize}
\item
What is the analogue of the Tate curve in the Calabi-Yau case?
Do we have some good model in this case like $(\widehat{\mathbb{G}_m}, \phi_{can})$?
\item 
In the case of the elliptic curve, it seems that the following is ture: the Ramanujan vector field $\sf{R}$ restricting to $t_1 = 0$ and further restricting to the Tate curve coincides with the pull back of the derivation $x\frac{\partial}{\partial x}$ of $\widehat{\mathbb{G}_m}$ under $\phi_{can}$.
\end{itemize}


















\end{document}



